\chapter{Bilangan Real} \label{rn:chapter}

%%%%%%%%%%%%%%%%%%%%%%%%%%%%%%%%%%%%%%%%%%%%%%%%%%%%%%%%%%%%%%%%%%%%%%%%%%%%%%

\section{Sifat-sifat dasar} \label{sec:basicpropsrn}

\sectionnotes{1,5 pertemuan kuliah}

Dalam analisis, objek utama yang kita kaji adalah himpunan
\myindex{bilangan real}.  Karena himpunan ini begitu mendasar, sering kali banyak
waktu dihabiskan untuk membangun himpunan bilangan real secara formal.  Namun, kita
akan mengambil pendekatan yang lebih mudah dan mengasumsikan bahwa ada suatu himpunan
dengan sifat-sifat yang tepat.
Tiga sifat utama bilangan real adalah bahwa himpunan ini
terurut, lengkap terhadap urutan tersebut,
dan merupakan lapangan yang serasi dengan urutan tersebut.
Kita mulai dengan urutan.

\begin{defn}
Suatu \emph{\myindex{himpunan terurut}} adalah himpunan $S$ yang dilengkapi
dengan relasi $<$ sedemikian sehingga
%\begin{enumerate}[(i),itemsep=0.5\itemsep,parsep=0.5\parsep,topsep=0.5\topsep,partopsep=0.5\partopsep]
%\begin{enumerate}[(i),nolistsep]
\begin{enumerate}[(i)]
\item \emph{(\myindex{trikotomi})} Untuk semua $x, y \in S$, tepat satu di antara
$x < y$, $x=y$, atau $y < x$ berlaku.
\item \emph{(transitivitas)} Jika $x,y,z \in S$ sedemikian sehingga $x < y$ dan $y
< z$, maka $x < z$.
\end{enumerate}
Kita menulis $x \leq y$ jika $x < y$ atau $x=y$.  Kita mendefinisikan
$>$ dan $\geq$ dengan cara yang wajar.
\glsadd{not:ordereq}
\glsadd{not:orderlt}
\glsadd{not:orderleq}
\glsadd{not:ordergt}
\glsadd{not:ordergeq}
\end{defn}


Himpunan bilangan rasional $\Q$ adalah himpunan terurut: Kita mengatakan
$x < y$ jika dan hanya jika $y-x$ adalah bilangan rasional positif, yaitu,
jika $y-x = \nicefrac{p}{q}$ dengan $p,q \in \N$.  Demikian pula,
$\N$ dan $\Z$ juga merupakan himpunan terurut.

Ada himpunan terurut yang bukan himpunan bilangan.  Sebagai contoh, himpunan
negara dapat diurutkan menurut luas daratannya, sehingga India $>$
Liechtenstein.
Salah satu himpunan terurut yang lazim dan telah Anda gunakan sejak sekolah dasar
adalah kamus.  Kamus merupakan himpunan kata terurut dengan urutan yang
disebut urutan leksikografis.  Himpunan terurut semacam itu sering muncul, misalnya,
dalam ilmu komputer.  Dalam buku ini, kita terutama akan mempelajari himpunan
bilangan terurut.

\begin{defn}
Misalkan $E \subset S$, dengan $S$ suatu himpunan terurut.
%\begin{enumerate}[(i),itemsep=0.5\itemsep,parsep=0.5\parsep,topsep=0.5\topsep,partopsep=0.5\partopsep]
%\begin{enumerate}[(i),nolistsep]
\begin{enumerate}[(i)]
\item Jika terdapat $b \in S$ sedemikian sehingga $x \leq b$ untuk semua $x \in E$,
maka $E$ disebut \emph{\myindex{terbatas di atas}} dan $b$
merupakan suatu \emph{\myindex{batas atas}} bagi $E$.
\item Jika terdapat $b \in S$ sedemikian sehingga $x \geq b$ untuk semua $x \in E$,
maka $E$ disebut \emph{\myindex{terbatas di bawah}} dan $b$
merupakan suatu \emph{\myindex{batas bawah}} bagi $E$.
\item Jika terdapat batas atas $b_0$ bagi $E$ sedemikian sehingga
$b_0 \leq b$ untuk setiap batas atas $b$ bagi $E$, maka
$b_0$ disebut \emph{\myindex{batas atas terkecil}} atau
\emph{\myindex{supremum}}
dari $E$.  Lihat \figureref{fig:lub}.  Kita menulis\glsadd{not:sup}
\begin{equation*}
\sup\, E \coloneqq b_0 .
\end{equation*}
\item Demikian pula, jika terdapat batas bawah $b_0$ bagi $E$ sedemikian sehingga
$b_0 \geq b$ untuk setiap batas bawah $b$ bagi $E$, maka
$b_0$ disebut \emph{\myindex{batas bawah terbesar}} atau
\emph{\myindex{infimum}}
dari $E$.  Kita menulis\glsadd{not:inf}
\begin{equation*}
\inf\, E \coloneqq b_0 .
\end{equation*}
\end{enumerate}
Jika suatu himpunan $E$ terbatas di atas sekaligus terbatas di bawah, kita cukup mengatakan bahwa
$E$ \emph{terbatas}\index{himpunan terbatas}.
\end{defn}

Notasi $\sup E$ dan $\inf E$ dapat dibenarkan karena
supremum (atau infimum) itu tunggal (jika ada): Jika $b$ dan
$b'$ keduanya merupakan supremum $E$, maka $b \leq b'$ dan $b' \leq b$, karena
$b$ dan $b'$ sama-sama merupakan batas atas terkecil, sehingga $b=b'$.


\begin{myfigureht}
\myincludepdft{lub}{%
Diagram himpunan E pada garis bilangan real yang digambarkan sebagai garis horizontal.
Sebuah panah di kiri menunjukkan bahwa bilangan semakin kecil ketika kita bergerak ke
kiri, sedangkan panah lain di kanan menunjukkan bahwa bilangan semakin besar
ketika kita bergerak ke kanan.  Himpunan E ditampilkan sebagai bagian-bagian garis
bilangan real yang disorot.  Beberapa tanda vertikal yang semuanya berada di sebelah kanan
himpunan E ditandai sebagai batas atas E, sedangkan tanda terkecil di antaranya,
yang tepat berada di tepi E, ditandai sebagai batas atas terkecil E.}
\caption{Himpunan $E$ yang terbatas di atas dan batas atas terkecil dari $E$.\label{fig:lub}}
\end{myfigureht}

Sebuah contoh sederhana:
Misalkan $S \coloneqq \{ a, b, c, d, e \}$ diurutkan sebagai $a < b < c < d < e$, dan
misalkan $E \coloneqq \{ a, c \}$.  Maka $c$, $d$, dan $e$ adalah batas atas bagi $E$, dan
$c$ adalah batas atas terkecil atau supremum dari $E$.

Supremum atau infimum bagi $E$ (sekalipun ada) tidak harus
berada di dalam $E$.  Himpunan $E \coloneqq \{ x \in \Q : x < 1 \}$ memiliki batas atas
terkecil 1, tetapi 1 tidak berada di dalam himpunan $E$ itu sendiri.
Himpunan
$G \coloneqq \{ x \in \Q : x \leq 1 \}$ juga memiliki batas atas 1,
dan dalam kasus ini $1 \in G$.
Himpunan $P \coloneqq \{ x \in \Q : x \geq 0 \}$ tidak memiliki batas atas (mengapa?) sehingga
tidak mungkin memiliki batas atas terkecil.  Himpunan $P$ memang memiliki batas bawah
terbesar, yaitu~0.

\begin{defn} \label{defn:lub}
Suatu himpunan terurut $S$ memiliki \emph{\myindex{sifat batas atas terkecil}} jika
setiap himpunan bagian tak kosong
$E \subset S$ yang terbatas di atas memiliki batas atas terkecil,
yaitu, $\sup\, E$ ada di dalam $S$.
\end{defn}

\emph{Sifat batas atas terkecil}
kadang-kadang disebut \emph{\myindex{sifat kelengkapan}} atau
\emph{\myindex{sifat kelengkapan Dedekind}}%
\footnote{%
Diambil dari nama matematikawan Jerman
\href{https://en.wikipedia.org/wiki/Richard_Dedekind}{Julius Wilhelm Richard Dedekind}
(1831--1916).}.
Bilangan real memiliki sifat ini.

\begin{example}
Himpunan bilangan rasional $\Q$ tidak memiliki sifat batas atas terkecil.  Himpunan bagian
$\{ x \in \Q : x^2 < 2 \}$ tidak memiliki supremum di $\Q$.  Nanti
kita akan melihat (\exampleref{example:sqrt2}) bahwa supremumnya adalah
$\sqrt{2}$, yang bukan bilangan rasional%
\footnote{Hal ini juga berlaku untuk semua akar lain dari 2, dan menariknya,
fakta bahwa $\sqrt[k]{2}$ tidak pernah rasional untuk $k > 1$ berkaitan dengan
kemustahilan menyetem piano
secara sempurna dalam semua tangga nada.  Lihat, misalnya:
\url{https://youtu.be/1Hqm0dYKUx4}.}.
Andaikan $x \in \Q$ sedemikian sehingga $x^2 = 2$.
Tuliskan $x=\nicefrac{m}{n}$ dalam bentuk paling sederhana.  Maka ${(\nicefrac{m}{n})}^2 = 2$
atau
$m^2 = 2n^2$.  Jadi, $m^2$ habis dibagi 2, sehingga $m$ habis dibagi
2.  Tuliskan $m = 2k$, sehingga ${(2k)}^2 = 2n^2$.  Bagilah dengan 2
dan perhatikan bahwa $2k^2 = n^2$, sehingga $n$ habis dibagi 2.  Namun, hal ini
bertentangan dengan fakta bahwa $\nicefrac{m}{n}$ telah ditulis dalam bentuk paling sederhana.
\end{example}

Ketiadaan sifat batas atas terkecil pada $\Q$ merupakan salah satu alasan terpenting
kita bekerja dengan $\R$ dalam analisis.  Himpunan $\Q$
sudah memadai bagi para ahli aljabar.  Namun, kami para analis memerlukan sifat batas atas
terkecil untuk dapat bekerja.
Kita juga mengharuskan bilangan real memiliki banyak sifat aljabar.
Khususnya, kita mengharuskan bilangan real membentuk suatu lapangan.


%\enlargethispage{\baselineskip}
\begin{defn}
Suatu himpunan $F$ disebut \emph{\myindex{lapangan}} jika padanya didefinisikan dua operasi,
yaitu penjumlahan $x+y$ dan perkalian $xy$, serta jika himpunan itu memenuhi
aksioma-aksioma\index{aksioma} berikut:
%mbxSTARTIGNORE
\begin{enumerate}[({A}1)]
%mbxENDIGNORE
%mbxCLOSEPARAGRAPH
%mbx <dl>
\item
%mbx <title>(A1)</title>
Jika $x \in F$ dan $y \in F$, maka $x+y \in F$.
\item
%mbx <title>(A2)</title>
\emph{(komutativitas penjumlahan)}
$x+y = y+x$ untuk semua $x,y \in F$.
\item
%mbx <title>(A3)</title>
\emph{(asosiativitas penjumlahan)}
$(x+y)+z = x+(y+z)$ untuk semua $x,y,z \in F$.
\item
%mbx <title>(A4)</title>
Terdapat elemen $0 \in F$ sedemikian sehingga
$0+x = x$ untuk semua $x \in F$.
\item
%mbx <title>(A5)</title>
Untuk setiap elemen $x\in F$, terdapat elemen $-x \in F$
sedemikian sehingga $x + (-x) = 0$.
%mbxCLOSEITEM
%mbx </dl>
%mbxSTARTIGNORE
\end{enumerate}
\begin{enumerate}[({M}1)]
%mbxENDIGNORE
%mbx <dl>
\item
%mbx <title>(M1)</title>
Jika $x \in F$ dan $y \in F$, maka $xy \in F$.
\item
%mbx <title>(M2)</title>
\emph{(komutativitas perkalian)}
$xy = yx$ untuk semua $x,y \in F$.
\item
%mbx <title>(M3)</title>
\emph{(asosiativitas perkalian)}
$(xy)z = x(yz)$ untuk semua $x,y,z \in F$.
\item
%mbx <title>(M4)</title>
Terdapat elemen $1 \in F$ (dan $1 \neq 0$) sedemikian sehingga
$1x = x$ untuk semua $x \in F$.
\item
%mbx <title>(M5)</title>
Untuk setiap $x\in F$ dengan $x \neq 0$, terdapat elemen
$\nicefrac{1}{x} \in F$
sedemikian sehingga $x(\nicefrac{1}{x}) = 1$.
%\end{enumerate}
%\begin{enumerate}
%mbxSTARTIGNORE
\item[(D)]
%mbxENDIGNORE
%mbxlatex \item
%mbx <title>(D)</title>
\emph{(hukum distributif)} $x(y+z) = xy+xz$
untuk semua $x,y,z \in F$.
%mbxSTARTIGNORE
\end{enumerate}
%mbxENDIGNORE
%mbxCLOSEITEM
%mbx </dl>
\end{defn}

\begin{example}
Himpunan bilangan rasional $\Q$ adalah lapangan.  Sebaliknya, $\Z$ bukan
lapangan karena tidak memuat invers perkalian.  Sebagai contoh,
tidak ada $x \in \Z$ sedemikian sehingga $2x = 1$, jadi (M5) tidak terpenuhi.  Anda
dapat memeriksa bahwa (M5) adalah satu-satunya sifat yang gagal\footnote{Seorang ahli aljabar akan mengatakan bahwa $\Z$ adalah
cincin terurut, atau mungkin cincin komutatif terurut.}.
\end{example}

Kita akan mengasumsikan fakta-fakta dasar tentang lapangan yang mudah
dibuktikan dari aksioma-aksioma tersebut.  Sebagai contoh, $0x = 0$ mudah dibuktikan
dengan memperhatikan bahwa $xx = (0+x)x = 0x+xx$, menggunakan (A4), (D)\@, dan (M2).  Kemudian,
dengan menggunakan (A5) pada $xx$, bersama (A2), (A3), dan (A4), kita memperoleh $0 = 0x$.

\begin{defn}
Suatu lapangan $F$ disebut \emph{\myindex{lapangan terurut}} jika
$F$ juga merupakan himpunan terurut sedemikian sehingga
\begin{enumerate}[(i)]
\item \label{defn:ordfield:i} Untuk $x,y,z \in F$, $x < y$ mengakibatkan $x+z <
y+z$.
\item \label{defn:ordfield:ii} Untuk $x,y \in F$, $x > 0$ dan $y > 0$
mengakibatkan $xy > 0$.
\end{enumerate}
Jika $x > 0$, kita mengatakan bahwa $x$ \emph{\myindex{positif}}.
Jika $x < 0$, kita mengatakan bahwa $x$ \emph{\myindex{negatif}}.
Kita juga mengatakan bahwa $x$ \emph{\myindex{nonnegatif}} jika $x \geq 0$,
dan $x$ \emph{\myindex{nonpositif}} jika $x \leq 0$.
\end{defn}

Bilangan rasional $\Q$ dengan urutan bakunya merupakan lapangan terurut.
Perinciannya kita serahkan kepada pembaca yang berminat.

\begin{prop} \label{prop:bordfield}
Misalkan $F$ suatu lapangan terurut dan $x,y,z,w \in F$.  Maka
\begin{enumerate}[(i)]
\item \label{prop:bordfield:i} Jika $x > 0$, maka $-x < 0$ (dan sebaliknya).
\item \label{prop:bordfield:ii} Jika $x > 0$ dan $y < z$, maka $xy < xz$.
\item \label{prop:bordfield:iii} Jika $x < 0$ dan $y < z$, maka $xy > xz$.
\item \label{prop:bordfield:iv} Jika $x \neq 0$, maka $x^2 > 0$.
\item \label{prop:bordfield:v} Jika $0 < x < y$, maka $0 < \nicefrac{1}{y} < \nicefrac{1}{x}$.
\item \label{prop:bordfield:vi} Jika $0 < x < y$, maka $x^2 < y^2$.
\item \label{prop:bordfield:vii} Jika $x \leq y$ dan $z \leq w$, maka $x + z \leq y + w$.
\end{enumerate}
\end{prop}

Perhatikan bahwa \ref{prop:bordfield:iv} secara khusus mengakibatkan $1 > 0$.

\begin{proof}
Mari kita buktikan \ref{prop:bordfield:i}.  Dari ketaksamaan $x > 0$ dan butir
\ref{defn:ordfield:i} dalam definisi lapangan terurut diperoleh
$x + (-x) > 0 + (-x)$.  Dengan menerapkan sifat-sifat aljabar lapangan, kita
memperoleh $0 > -x$.  Arah \myquote{sebaliknya} diperoleh melalui perhitungan serupa.

Untuk \ref{prop:bordfield:ii}, perhatikan bahwa $y < z$ mengakibatkan
$0 < z - y$ menurut
butir \ref{defn:ordfield:i} dalam definisi lapangan terurut.
Terapkan butir
\ref{defn:ordfield:ii} dalam definisi lapangan terurut untuk memperoleh
$0 < x(z-y)$.  Berdasarkan sifat-sifat aljabar, $0 < xz - xy$.
Dengan kembali menggunakan butir
\ref{defn:ordfield:i} dalam definisi tersebut, diperoleh $xy < xz$.

Bagian \ref{prop:bordfield:iii} diserahkan sebagai latihan.

Untuk membuktikan bagian \ref{prop:bordfield:iv}, pertama andaikan $x > 0$.
Menurut butir \ref{defn:ordfield:ii} dalam definisi lapangan terurut,
$x^2 > 0$ (gunakan $y=x$).  Jika $x < 0$, kita menggunakan
bagian \ref{prop:bordfield:iii} dari proposisi ini dengan mengambil $y=x$ dan
$z=0$.

Untuk membuktikan bagian \ref{prop:bordfield:v}, perhatikan bahwa
$\nicefrac{1}{x}$ tidak mungkin sama dengan nol (mengapa?).
Andaikan $\nicefrac{1}{x} < 0$,
maka $\nicefrac{-1}{x} > 0$ menurut \ref{prop:bordfield:i}.
Terapkan butir \ref{defn:ordfield:ii} dalam definisi (karena $x > 0$) untuk memperoleh
$x(\nicefrac{-1}{x}) > 0$ atau $-1 > 0$, yang bertentangan dengan $1 > 0$ setelah kembali menggunakan bagian
\ref{prop:bordfield:i}.  Jadi $\nicefrac{1}{x} > 0$.
Demikian pula, $\nicefrac{1}{y} > 0$.  Dengan demikian $(\nicefrac{1}{x})(\nicefrac{1}{y}) > 0$
menurut definisi lapangan terurut, dan berdasarkan bagian \ref{prop:bordfield:ii},
\begin{equation*}
(\nicefrac{1}{x})(\nicefrac{1}{y})x < (\nicefrac{1}{x})(\nicefrac{1}{y})y .
\end{equation*}
Berdasarkan sifat-sifat aljabar, $\nicefrac{1}{y} < \nicefrac{1}{x}$.

Bagian \ref{prop:bordfield:vi} dan \ref{prop:bordfield:vii} diserahkan sebagai
latihan.
\end{proof}

Hasil kali dua bilangan positif (elemen suatu lapangan terurut) bernilai positif
menurut definisi lapangan terurut.
Namun, tidak benar bahwa jika hasil kalinya positif, maka masing-masing dari kedua
faktornya harus positif.  Sebagai contoh, $(-1)(-1) = 1 > 0$.

\begin{prop}
Misalkan $x,y \in F$, dengan $F$ suatu lapangan terurut.  Jika
$xy > 0$, maka $x$ dan $y$ keduanya positif atau keduanya negatif.
\end{prop}

\begin{proof}
Kita buktikan kontraposisinya: \emph{Jika $x$ atau $y$ bernilai nol, atau
jika $x$ dan $y$ berlainan tanda, maka $xy$ tidak positif.}
Jika $x$ atau $y$ bernilai nol, maka $xy$ bernilai nol sehingga tidak positif.
Jadi, andaikan $x$ dan $y$ tak nol dan
berlainan tanda.
Tanpa mengurangi keumuman,
andaikan $x > 0$ dan $y < 0$.
Kalikan $y < 0$ dengan $x$ untuk memperoleh $xy < 0x = 0$.
\end{proof}

\begin{example} \label{example:complexfield}
Pembaca mungkin pernah mendengar tentang \emph{\myindex{bilangan kompleks}},
yang biasanya dilambangkan dengan $\C$.\glsadd{not:complex}
Artinya, $\C$ adalah himpunan bilangan berbentuk
$x + iy$, dengan $x$ dan $y$ bilangan real, dan $i$ adalah
bilangan imajiner, yakni bilangan sedemikian sehingga $i^2 = -1$.  Pembaca mungkin
ingat dari aljabar bahwa $\C$ juga merupakan lapangan; namun, himpunan ini bukan
lapangan terurut.  Walaupun $\C$ dapat dijadikan himpunan terurut dengan suatu cara,
tidak mungkin memberikan
urutan pada $\C$ yang menjadikannya lapangan terurut: Dalam setiap lapangan terurut,
$-1 < 0$ dan $x^2 > 0$ untuk setiap $x$ yang tak nol, tetapi dalam $\C$, $i^2 = -1$.
\end{example}

Terakhir, suatu lapangan terurut yang memiliki sifat batas atas terkecil juga memiliki
sifat yang bersesuaian untuk batas bawah terbesar.

\begin{prop}
Misalkan $F$ suatu lapangan terurut dengan sifat batas atas terkecil.
Misalkan $A \subset F$ suatu himpunan tak kosong yang terbatas di bawah.
Maka $\inf\, A$ ada.
\end{prop}

\begin{proof}
Misalkan $B \coloneqq \{ -x : x \in A \}$. Misalkan $b \in F$ suatu batas bawah bagi $A$:
Jika $x \in A$, maka $x \geq b$
sehingga $-x \leq -b$.  Jadi $-b$
merupakan batas atas bagi $B$.
Karena $F$ memiliki sifat batas atas terkecil, $c\coloneqq\sup\, B$ ada, dan $c \leq -b$.
Karena $y \leq c$
untuk semua $y \in B$, kita memperoleh $-c \leq x$ untuk semua $x \in A$.
Jadi,
$-c$ merupakan batas bawah bagi $A$.  Karena $-c \geq b$
dan $b$ adalah sebarang batas bawah bagi $A$,
batas bawah terbesar dari $A$ ada dan sama dengan $-c$.
\end{proof}

\subsection{Latihan}

\begin{exercise}
Buktikan bagian \ref{prop:bordfield:iii} dari \propref{prop:bordfield}.
Dengan kata lain, misalkan $F$ suatu lapangan terurut dan $x,y,z \in F$.  Buktikan
bahwa
jika $x < 0$ dan $y < z$, maka $xy > xz$.
\end{exercise}


\begin{exercise} \label{exercise:finitesethasminmax}
Misalkan $S$ suatu himpunan terurut.  Misalkan $A \subset S$ suatu himpunan bagian hingga tak kosong.
Tunjukkan bahwa $A$ terbatas.  Lebih lanjut,
tunjukkan bahwa
$\inf\, A$ ada dan berada di $A$, serta
$\sup\, A$ ada dan berada di $A$.
Petunjuk: Gunakan \hyperref[induction:thm]{induksi}.
\end{exercise}

\begin{exercise} \label{exercise:squareineq}
Buktikan bagian \ref{prop:bordfield:vi} dari \propref{prop:bordfield}.
Dengan kata lain, misalkan $x, y \in F$, dengan $F$ suatu lapangan terurut, sedemikian sehingga
$0 < x < y$.  Tunjukkan bahwa $x^2 < y^2$.
\end{exercise}

\begin{exercise}
Misalkan $S$ suatu himpunan terurut.  Misalkan $B \subset S$ terbatas (di atas dan
di bawah).  Misalkan $A \subset B$ suatu himpunan bagian tak kosong.
Andaikan semua $\inf$ dan
$\sup$ tersebut ada. Tunjukkan bahwa
\begin{equation*}
\inf\, B \leq \inf\, A \leq \sup\, A \leq \sup\, B .
\end{equation*}
\end{exercise}

\begin{exercise}
Misalkan $S$ suatu himpunan terurut.  Misalkan $A \subset S$ dan andaikan
$b$ merupakan batas atas bagi $A$.  Andaikan $b \in A$.  Tunjukkan
bahwa $b = \sup\, A$.
\end{exercise}

\begin{exercise}
Misalkan $S$ suatu himpunan terurut.  Misalkan $A \subset S$ tak kosong dan
terbatas di atas.
Andaikan $\sup\, A$ ada dan
$\sup\, A \notin A$.
Tunjukkan bahwa $A$ memuat suatu himpunan bagian
terhitung tak hingga.
%In particular, $A$ is infinite.
\end{exercise}

\begin{exercise}
Temukan suatu urutan tak baku pada himpunan bilangan asli $\N$
sedemikian sehingga terdapat himpunan bagian sejati tak kosong $A \subsetneq \N$
dan sedemikian sehingga $\sup\, A$ ada di $\N$, tetapi $\sup\, A \notin A$.
Agar tidak membingungkan, sebaiknya gunakan notasi lain
untuk urutan tak baku tersebut, misalnya $n \prec m$.
\end{exercise}

\begin{exercise}
\pagebreak[2]
Misalkan $F \coloneqq \{ 0, 1, 2 \}$.
\begin{enumerate}[a)]
\item
Buktikan bahwa hanya ada
satu cara untuk mendefinisikan penjumlahan dan perkalian yang menjadikan $F$ suatu
lapangan, dengan $0$ dan $1$ mempunyai arti lazim seperti pada (A4) dan (M4).
\item
Tunjukkan bahwa $F$ tidak mungkin merupakan lapangan terurut.
\end{enumerate}
\end{exercise}

\begin{exercise} \label{exercise:dominatingb}
\pagebreak[2]
Misalkan $S$ suatu himpunan terurut dan
misalkan $A$ suatu himpunan bagian tak kosong sedemikian sehingga $\sup \, A$ ada.  Andaikan terdapat
$B \subset A$ sedemikian sehingga untuk setiap $x \in A$ terdapat $y \in B$
sedemikian sehingga $x \leq y$.  Tunjukkan bahwa $\sup \, B$ ada dan $\sup \, B = \sup \, A$.
\end{exercise}

\begin{exercise}
Misalkan $D$ himpunan terurut yang terdiri atas semua kata yang mungkin (bukan hanya kata-kata bahasa Inggris,
melainkan semua untai huruf dengan panjang sebarang)
yang menggunakan alfabet Latin dan hanya huruf kecil.  Urutannya adalah
urutan leksikografis seperti dalam kamus (mis.\ \emph{aa} $<$ \emph{aaa} $<$ \emph{dog} $<$ \emph{door}).
Misalkan $A$ himpunan bagian dari $D$ yang memuat kata-kata dengan
huruf pertama `a' (mis.\ \emph{a} $\in A$, \emph{abcd} $\in A$).
Tunjukkan bahwa $A$ memiliki supremum dan tentukan supremum tersebut.
\end{exercise}

\begin{exercise}
Misalkan $F$ suatu lapangan terurut dan $x,y,z,w \in F$.
\begin{enumerate}[a)]
\item
Buktikan bagian \ref{prop:bordfield:vii} dari \propref{prop:bordfield}.
Dengan kata lain,
jika $x \leq y$ dan $z \leq w$, maka $x+z \leq y+w$.
\item
Buktikan bahwa
jika $x < y$ dan $z \leq w$, maka $x+z < y+w$.
\end{enumerate}
\end{exercise}

\begin{exercise}
Buktikan bahwa setiap lapangan terurut harus memuat suatu himpunan terhitung tak hingga.
\end{exercise}

\begin{exercise}
Misalkan $\N_{\infty} \coloneqq \N \cup \{ \infty \}$, dengan unsur-unsur $\N$
mengikuti urutan biasa di antara sesamanya,
dan $k < \infty$ untuk setiap $k \in \N$.  Tunjukkan bahwa $\N_{\infty}$ merupakan himpunan terurut
dan bahwa setiap himpunan bagian $E \subset \N_{\infty}$ memiliki supremum di
$\N_{\infty}$ (pastikan juga menangani kasus himpunan kosong).
Catatan: Perhatikan bahwa kasus himpunan kosong akan berbeda
jika himpunan terurut kita adalah bilangan real atau bilangan
real diperluas dari bagian berikutnya.  Penting bahwa
himpunan terurut kita di sini adalah $\N_{\infty}$.
\end{exercise}

\begin{exercise}
Misalkan $S \coloneqq \{ a_k : k \in \N \} \cup \{ b_k : k \in \N \}$, diurutkan sedemikian sehingga
$a_k < b_j$ untuk setiap $k$ dan $j$,
$a_k < a_m$ setiap kali $k < m$,
dan $b_k > b_m$ setiap kali $k < m$.
\begin{enumerate}[a)]
\item
Tunjukkan bahwa $S$ merupakan himpunan terurut.
\item
Tunjukkan bahwa setiap himpunan bagian dari $S$ terbatas (baik di atas maupun di bawah).
\item
Temukan himpunan bagian terbatas dari $S$ yang tidak memiliki batas atas terkecil.
\end{enumerate}
\end{exercise}

%%%%%%%%%%%%%%%%%%%%%%%%%%%%%%%%%%%%%%%%%%%%%%%%%%%%%%%%%%%%%%%%%%%%%%%%%%%%%%

\sectionnewpage
\section{Himpunan bilangan real} \label{sec:setofreals}

%mbxINTROSUBSECTION

\sectionnotes{2 pertemuan kuliah; bilangan real diperluas bersifat opsional}

\subsection{Himpunan bilangan real}

Akhirnya kita sampai pada sistem bilangan real.  Untuk menyederhanakan pembahasan, alih-alih mengonstruksi
himpunan bilangan real dari bilangan rasional, kita cukup menyatakan keberadaannya
sebagai teorema tanpa bukti.  Perhatikan bahwa $\Q$ merupakan
lapangan terurut.

\begin{thm}
Terdapat suatu lapangan terurut unik\footnote{Keunikan dipahami hingga isomorfisme, tetapi kita ingin
menghindari penggunaan aljabar yang berlebihan.  Bagi keperluan kita, cukup mengasumsikan
bahwa suatu himpunan bilangan real ada.  Lihat Rudin~\cite{Rudin:baby} untuk
konstruksi dan rincian lebih lanjut.}
$\R$ yang memiliki \hyperref[defn:lub]{sifat batas atas terkecil}
sedemikian sehingga $\Q \subset \R$.
\end{thm}

Perhatikan pula bahwa $\N \subset \Q$.  Kita telah melihat bahwa $1 > 0$.
Dengan \hyperref[induction:thm]{induksi} (latihan),
kita dapat membuktikan bahwa $n > 0$ untuk setiap $n \in \N$.
Dengan cara serupa, kita memverifikasi pernyataan-pernyataan sederhana tentang bilangan rasional.
Sebagai contoh, kita telah membuktikan bahwa jika $n > 0$, maka $\nicefrac{1}{n} > 0$.
Selanjutnya, $m < k$ mengakibatkan $\nicefrac{m}{n} < \nicefrac{k}{n}$.

Analisis berkaitan dengan pembuktian ketaksamaan, dan
proposisi berikut, atau salah satu dari banyak variasinya,
merupakan cara seorang analis membuktikan ketaksamaan tak ketat.

\begin{prop}
Jika $x \in \R$ sedemikian sehingga $x \leq \epsilon$ untuk setiap
$\epsilon \in \R$ dengan
$\epsilon > 0$, maka $x \leq 0$.
\end{prop}

\begin{proof}
Jika $x > 0$, maka $0 < \nicefrac{x}{2} < x$ (mengapa?).  Ambil
$\epsilon = \nicefrac{x}{2}$ sehingga diperoleh kontradiksi.  Jadi $x \leq 0$.
\end{proof}

Untuk $x$ nonnegatif, diperoleh kesamaan:
\emph{Jika $x \geq 0$ sedemikian sehingga $x \leq \epsilon$ untuk setiap
$\epsilon > 0$, maka $x = 0$.}
Versi yang umum menggunakan nilai mutlak (lihat \sectionref{sec:absval}):
\emph{Jika $\sabs{x} \leq \epsilon$ untuk setiap
$\epsilon > 0$, maka $x = 0$.}
Untuk membuktikan $x \geq 0$, seorang analis dapat membuktikan
bahwa $x \geq -\epsilon$ untuk setiap $\epsilon > 0$.
Mulai sekarang, ketika kita menulis $x \geq 0$ atau
$\epsilon > 0$, secara otomatis kita maksudkan bahwa $x \in \R$ dan $\epsilon \in \R$.

Gagasan di balik proposisi di atas adalah bahwa
setiap kali kita memiliki dua bilangan real $a < b$, terdapat bilangan
real lain $c$ sedemikian sehingga
$a < c < b$.  Terdapat tak hingga banyaknya $c$ seperti itu.  Salah satunya, misalnya, adalah
$c = \frac{a+b}{2}$ (mengapa?).
Kita akan menggunakan fakta ini dalam contoh berikutnya.

Hal yang paling berguna dari $\R$ bagi para analis
bukan sekadar sifatnya sebagai lapangan terurut, melainkan bahwa lapangan ini memiliki
\hyperref[defn:lub]{sifat batas atas terkecil}.
Pada dasarnya, kita menginginkan $\Q$, tetapi kita juga
ingin dapat mengambil supremum (dan infimum) sesuka hati.
Jadi, kita mengambil $\Q$ dan menambahkan cukup banyak bilangan untuk memperoleh $\R$.

Kita telah menyebutkan bahwa $\R$ memuat unsur-unsur yang tidak berada di $\Q$
karena \hyperref[defn:lub]{sifat batas atas terkecil}.  Mari kita
buktikan hal itu.
Kita telah melihat bahwa tidak ada
akar kuadrat rasional dari dua.  Himpunan
$\{ x \in \Q : x^2 < 2 \}$ mengakibatkan keberadaan bilangan real
$\sqrt{2}$, meskipun fakta ini memerlukan sedikit usaha.  Lihat pula
\exerciseref{exercise:sqrt2QorR}.

\begin{example} \label{example:sqrt2}
Klaim: \emph{Terdapat suatu bilangan positif unik
$r \in \R$ sedemikian sehingga $r^2 = 2$.  Kita melambangkan $r$ dengan $\sqrt{2}$.}

\begin{proof}
Ambil himpunan
$A \coloneqq \{ x \in \R : x^2 < 2 \}$.  Pertama-tama kita tunjukkan bahwa $A$ terbatas di atas dan
tak kosong.  Ketaksamaan $x \geq 2$ mengakibatkan $x^2 \geq 4$
(lihat \exerciseref{exercise:squareineq}).
Jadi, jika $x^2 < 2$, maka $x < 2$.  Dengan demikian, $A$ terbatas di atas.
Karena $1 \in A$, himpunan $A$ tak kosong.  Oleh karena itu, supremumnya ada.

Misalkan $r \coloneqq \sup\, A$.  Kita akan menunjukkan bahwa $r^2 = 2$ dengan menunjukkan
bahwa $r^2 \geq 2$ dan $r^2 \leq 2$.  Beginilah cara para analis menunjukkan
kesamaan, yakni dengan menunjukkan dua ketaksamaan.
Kita telah mengetahui bahwa $r \geq 1 > 0$.

%\medskip
Pada uraian berikut, mungkin tampak seakan-akan beberapa bentuk tertentu muncul
begitu saja.  Ketika menulis bukti seperti ini, tentu saja kita baru menemukan
bentuk-bentuk tersebut setelah mencoba-coba berdasarkan apa yang ingin kita buktikan.  Urutan
penulisan bukti tidak harus sama dengan urutan ketika kita
menemukan bukti tersebut.

Pertama-tama mari kita tunjukkan bahwa $r^2 \geq 2$.
Ambil bilangan positif $s$ sedemikian sehingga $s^2 < 2$.  Kita ingin mencari $h > 0$
sedemikian sehingga ${(s+h)}^2 < 2$.
Karena $2-s^2 > 0$, kita memperoleh $\frac{2-s^2}{2s+1} > 0$.
Pilih $h \in \R$ sedemikian sehingga
$0 < h < \frac{2-s^2}{2s+1}$.
Selain itu, andaikan $h < 1$.  Kita peroleh taksiran:
\begin{equation*}
\begin{aligned}
{(s+h)}^2 - s^2 & = h(2s + h) \\
 & < h(2s+1) & & \quad \bigl(\text{karena } h < 1\bigr) \\
 & < 2-s^2 & & \quad \bigl(\text{karena } h < \tfrac{2-s^2}{2s+1} \bigr) .
\end{aligned}
\end{equation*}
Oleh karena itu, ${(s+h)}^2 < 2$.  Jadi $s+h \in A$, tetapi karena $h > 0$,
kita memperoleh $s+h > s$.  Maka $s < r = \sup\, A$.  Karena $s$ adalah sebarang
bilangan positif sedemikian sehingga $s^2 < 2$, diperoleh $r^2 \geq 2$.

%\medskip

Sekarang ambil bilangan positif $s$ sedemikian sehingga
$s^2 > 2$.  Kita ingin mencari $h > 0$ sedemikian sehingga
${(s-h)}^2 > 2$ dan $s-h$ tetap positif.
Karena $s^2-2 > 0$, kita memperoleh $\frac{s^2-2}{2s} > 0$.
Misalkan $h \coloneqq \frac{s^2-2}{2s}$,
dan periksa bahwa $s-h=s-\frac{s^2-2}{2s} = \frac{s}{2}+\frac{1}{s} > 0$.
Kita peroleh taksiran:
\begin{equation*}
\begin{aligned}
s^2 - {(s-h)}^2 & = 2sh - h^2 \\
 & < 2sh & & \quad \bigl( \text{karena } h^2 > 0 \text{ sebab } h \neq 0 \bigr) \\
 & = s^2-2 & & \quad \bigl( \text{karena } h = \tfrac{s^2-2}{2s} \bigr) .
\end{aligned}
\end{equation*}
Dengan mengurangkan $s^2$ dari kedua ruas lalu mengalikan kedua ruas dengan $-1$, kita memperoleh
${(s-h)}^2 > 2$.  Oleh karena itu, $s-h \notin A$.

Lebih lanjut, jika $x \geq s-h$,
maka $x^2 \geq {(s-h)}^2 > 2$ (karena $x > 0$ dan $s-h > 0$), sehingga $x \notin A$.
Jadi,
$s-h$ merupakan batas atas bagi $A$.  Namun, $s-h < s$, atau dengan kata
lain, $s > r = \sup\, A$.  Dengan demikian, $r^2 \leq 2$.

\medskip

Kedua ketaksamaan $r^2 \geq 2$ dan $r^2 \leq 2$ bersama-sama mengakibatkan
$r^2 = 2$.  Bagian pembuktian keberadaan telah selesai.  Kita masih perlu
membuktikan keunikan.  Andaikan $s \in \R$ sedemikian sehingga $s^2 = 2$ dan $s > 0$.
Jadi, $s^2 = r^2$.  Namun, jika $0 < s < r$, maka $s^2 < r^2$.  Dengan cara serupa,
$0 < r < s$ mengakibatkan $r^2 < s^2$.  Dengan demikian, $s = r$.
\end{proof}
\end{example}

Bilangan $\sqrt{2} \notin \Q$.  Himpunan
$\R \setminus \Q$ disebut himpunan bilangan
\emph{\myindex{irasional}}.  Kita baru saja membuktikan bahwa $\R \setminus \Q$
tak kosong.  Seperti akan kita lihat, himpunan ini bukan sekadar tak kosong,
melainkan benar-benar sangat besar.

Dengan menggunakan teknik yang sama seperti di atas, kita dapat menunjukkan bahwa bilangan real positif
$x^{1/n}$ ada untuk setiap $n\in \N$ dan setiap $x > 0$.
Artinya, untuk setiap $x > 0$,
terdapat bilangan real positif unik $r$ sedemikian sehingga $r^n = x$.
Buktinya diserahkan sebagai latihan.

\subsection{Sifat Archimedes}

Seperti telah kita lihat, terdapat banyak sekali bilangan real dalam setiap interval.  Namun,
terdapat pula tak hingga banyaknya bilangan rasional dalam setiap interval.  Pernyataan
berikut merupakan salah satu fakta mendasar tentang bilangan real.
Kedua bagian teorema berikut sebenarnya ekuivalen, meskipun mungkin
tidak tampak demikian pada pandangan pertama.


\begin{thm} \label{thm:arch}
%FIXMEevillayouthack
\pagebreak[3]
\leavevmode
\begin{enumerate}[(i)]
\item \label{thm:arch:i} \emph{(\myindex{Sifat Archimedes})}%
\footnote{Dinamai menurut matematikawan Yunani Kuno
\href{https://en.wikipedia.org/wiki/Archimedes}{Archimedes dari Sirakusa}
(sekitar\ 287 SM -- sekitar\ 212 SM)\@.
Sifat ini adalah Aksioma V dari karya Archimedes \myquote{Tentang Bola dan
Tabung} (225 SM)\@.}
Jika $x, y \in \R$ dan
$x > 0$, maka terdapat $n \in \N$ sedemikian sehingga
\begin{equation*}
nx > y .
\end{equation*}
\item \label{thm:arch:ii} \emph{($\Q$ rapat di $\R$\index{kepadatan
bilangan rasional})} Jika $x, y \in \R$ dan
$x < y$, maka terdapat $r \in \Q$ sedemikian sehingga
$x < r < y$.
\end{enumerate}
\end{thm}

\begin{proof}
Mari kita buktikan \ref{thm:arch:i}.  Bagi kedua ruas dengan $x$.
Maka \ref{thm:arch:i} menyatakan bahwa untuk setiap bilangan real $t\coloneqq \nicefrac{y}{x}$,
kita dapat menemukan $n \in \N$ sedemikian sehingga $n > t$.  Dengan kata lain,
\ref{thm:arch:i} menyatakan bahwa $\N \subset \R$ tidak terbatas di atas.
Andaikan, demi memperoleh kontradiksi, bahwa $\N$ terbatas di atas.  Misalkan $b \coloneqq \sup \N$.
Bilangan $b-1$ tidak mungkin merupakan batas atas bagi $\N$ karena bilangan itu lebih kecil
daripada $b$ (batas atas terkecil).  Jadi, terdapat $m \in \N$ sedemikian sehingga $m > b-1$.
Tambahkan satu untuk memperoleh $m+1 > b$, yang bertentangan dengan fakta bahwa $b$ merupakan
batas atas.

\begin{myfigureht}
\myincludepdft{figdensofQ}{%
Diagram garis bilangan real horizontal dengan beberapa tanda vertikal tinggi
yang berjarak sama, dengan jarak antartanda
sebesar 1 per n.  Tanda kedua
dari kiri diberi label m dikurangi 1, lalu seluruhnya dibagi n.
Setelah tanda ini terdapat bilangan x yang ditandai dengan tanda pendek.
Tanda tinggi berikutnya diberi label m dibagi n.
Kemudian terdapat bilangan y, lalu tanda tinggi berlabel
m ditambah 1, lalu seluruhnya dibagi n.  Setelahnya tampak dua tanda tinggi lagi.}
\caption{Gagasan bukti kerapatan $\Q$: Carilah $n$ sedemikian sehingga $y-x >
\nicefrac{1}{n}$, lalu ambil $m$ terkecil sedemikian sehingga $\nicefrac{m}{n} > x$.\label{figdensofQ}}
\end{myfigureht}
Mari kita tangani \ref{thm:arch:ii}.
Lihat \figureref{figdensofQ}
untuk ilustrasi gagasan di balik bukti tersebut.
Pertama, andaikan $x \geq 0$.
Perhatikan bahwa $y-x > 0$.
Berdasarkan \ref{thm:arch:i}, terdapat $n \in \N$ sedemikian sehingga
\begin{equation*}
n(y-x) > 1
\qquad \text{atau} \qquad
y-x > \nicefrac{1}{n}.
\end{equation*}
Sekali lagi berdasarkan \ref{thm:arch:i}, himpunan
$A \coloneqq \{ k \in \N : k > nx \}$ tak kosong.  Berdasarkan
\hyperlink{wop:link}{sifat keterurutan baik}
dari $\N$, $A$ memiliki unsur terkecil $m$.
Karena $m \in A$, kita memperoleh $m > nx$.
Bagi kedua ruas dengan $n$ untuk memperoleh $x < \nicefrac{m}{n}$.
Karena $m$ merupakan unsur terkecil
dari $A$, maka $m-1 \notin A$.
Jika $m > 1$, maka $m-1 \in \N$, tetapi $m-1 \notin A$ sehingga $m-1 \leq nx$.
Jika $m=1$,
maka $m-1 = 0$, dan $m-1 \leq nx$ tetap berlaku karena $x \geq 0$.
Dengan kata lain,
\begin{equation*}
m-1 \leq nx \qquad \text{atau} \qquad m \leq nx+1 .
\end{equation*}
Di sisi lain,
dari $n(y-x) > 1$ kita memperoleh $ny > 1+nx$.
Jadi, $ny > 1+nx \geq m$, dan karena itu $y > \nicefrac{m}{n}$.
Dengan menggabungkan semuanya, kita memperoleh $x < \nicefrac{m}{n} < y$.
Jadi, ambil $r = \nicefrac{m}{n}$.

Sekarang andaikan $x < 0$.  Jika $y > 0$, cukup ambil $r=0$.  Jika
$y \leq 0$, maka $0 \leq -y < -x$, dan kita
menemukan bilangan rasional $q$ sedemikian sehingga $-y < q < -x$.  Kemudian, ambil $r = -q$.
\end{proof}

Mari kita nyatakan dan buktikan suatu akibat sederhana tetapi berguna dari
\hyperref[thm:arch:i]{sifat Archimedes}.

\begin{cor}
$\inf \{ \nicefrac{1}{n} : n \in \N \} = 0$.
Lihat \figureref{fig:oneovernset}.
\end{cor}

\begin{proof}
Misalkan $A \coloneqq \{ \nicefrac{1}{n} : n \in \N \}$.  Jelas bahwa $A$ tak kosong.
Selain itu,
$\nicefrac{1}{n} > 0$ untuk setiap $n \in \N$, sehingga $0$ merupakan batas bawah
dan $b \coloneqq \inf\, A$ ada.
Karena $0$ merupakan batas bawah, $b \geq 0$.
Ambil sebarang $a > 0$.  Berdasarkan
\hyperref[thm:arch:i]{sifat Archimedes}, terdapat $n$ sedemikian sehingga
$na > 1$, yaitu $a > \nicefrac{1}{n} \in A$.  Oleh karena itu,
$a$ tidak mungkin merupakan batas bawah bagi $A$.  Dengan demikian, $b=0$.
\end{proof}

\begin{myfigureht}
\myincludegraphics{oneovernset}{%
Diagram garis bilangan real horizontal dengan tanda-tanda vertikal yang menandai
bilangan 1, 1 per 2, 1 per 3, 1 per 4, dan seterusnya.  Tanda-tanda ini semakin
mendekati 0, yang juga ditandai, dan tanda-tanda itu begitu rapat
ketika mencapai 0 sehingga di sana hanya tampak persegi panjang hitam.}
\caption{Himpunan $\{ \nicefrac{1}{n} : n \in \N \}$ dan infimumnya $0$.\label{fig:oneovernset}}
\end{myfigureht}



\subsection{Menggunakan supremum dan infimum}

Supremum dan infimum
kompatibel dengan operasi aljabar.  Untuk suatu himpunan $A \subset \R$ dan
$x \in \R$, definisikan
\begin{align*}
x + A & \coloneqq \{ x+y \in \R : y \in A \} , \\
xA & \coloneqq \{ xy \in \R : y \in A \} .
\end{align*}
Sebagai contoh, jika $A = \{ 1,2,3 \}$, maka $5+A = \{ 6,7,8 \}$ dan $3A = \{ 3,6,9
\}$.

\begin{prop} \label{prop:supinfalg}
Misalkan $A \subset \R$ tak kosong.
\begin{enumerate}[(i)]
\item Jika $x \in \R$ dan $A$ terbatas di atas, maka $\sup (x+A) = x + \sup\, A$.
\item Jika $x \in \R$ dan $A$ terbatas di bawah, maka $\inf (x+A) = x + \inf\, A$.
\item Jika $x > 0$ dan $A$ terbatas di atas, maka $\sup (xA) = x ( \sup\, A )$.
\item Jika $x > 0$ dan $A$ terbatas di bawah, maka $\inf (xA) = x ( \inf\, A )$.
\item Jika $x < 0$ dan $A$ terbatas di bawah, maka $\sup (xA) = x ( \inf\, A )$.
\item Jika $x < 0$ dan $A$ terbatas di atas, maka $\inf (xA) = x ( \sup\, A )$.
\end{enumerate}
\end{prop}

Perhatikan bahwa mengalikan suatu himpunan dengan bilangan negatif mempertukarkan supremum dan infimum.
Selain itu, karena proposisi tersebut mengakibatkan adanya
supremum (atau infimum,\ sesuai kasusnya) dari $x+A$ atau $xA$,
proposisi tersebut juga mengakibatkan bahwa $x+A$ atau $xA$
tak kosong dan terbatas di atas (atau di bawah,\ sesuai kasusnya).

\begin{proof}
Mari kita buktikan pernyataan pertama saja.  Sisanya diserahkan sebagai latihan.

Andaikan $b$ merupakan batas atas bagi $A$.  Artinya, $y \leq b$ untuk setiap $y \in A$.
Maka $x+y \leq x+b$ untuk setiap $y \in A$, sehingga $x+b$ merupakan batas
atas bagi $x+A$.  Secara khusus, jika $b = \sup\, A$, maka
\begin{equation*}
\sup (x+A) \leq x+b = x+ \sup\, A .
\end{equation*}

Ketaksamaan sebaliknya dibuktikan dengan cara serupa. Jika $c$ merupakan batas atas bagi $x+A$,
maka $x+y \leq c$
untuk setiap $y \in A$, sehingga $y \leq c-x$ untuk setiap $y \in A$.
Jadi, $c-x$ merupakan batas atas bagi $A$.
Jika $c = \sup (x+A)$, maka
\begin{equation*}
\sup\, A \leq c-x = \sup (x+A) -x .
\end{equation*}
Dengan demikian, hasilnya terbukti.
\end{proof}

Terkadang kita perlu menerapkan supremum atau infimum dua kali.  Berikut sebuah contoh.

\begin{prop} \label{infsupineq:prop}
Misalkan $A, B \subset \R$ himpunan-himpunan tak kosong sedemikian sehingga $x \leq y$ setiap kali $x \in A$ dan
$y \in B$.  Maka $A$ terbatas di atas, $B$ terbatas di bawah, dan $\sup\, A \leq \inf\, B$.
\end{prop}

\begin{proof}
Setiap $x \in A$ merupakan batas bawah bagi $B$.
Oleh karena itu, $x \leq \inf\, B$ untuk setiap $x \in A$,
sehingga $\inf\, B$ merupakan batas atas bagi $A$.
Dengan demikian, $\sup\, A \leq \inf\, B$.
\end{proof}

Kita harus berhati-hati terhadap ketaksamaan ketat ketika mengambil supremum dan infimum.
Perhatikan bahwa $x < y$ setiap kali $x \in A$ dan
$y \in B$ tetap hanya mengakibatkan $\sup\, A \leq \inf\, B$, bukan ketaksamaan
ketat.
Sebagai contoh, ambil $A \coloneqq \{ 0 \}$ dan $B \coloneqq \{ \nicefrac{1}{n}
: n \in \N \}$.
Maka $0 < \nicefrac{1}{n}$
untuk setiap $n \in \N$.  Namun, $\sup\, A = 0$ dan $\inf\, B = 0$.
Nuansa penting ini sering muncul.

\medskip
\pagebreak[1]

Bukti fakta berikut, yang
sering digunakan, diserahkan kepada pembaca.
Hasil serupa berlaku untuk infimum.

\begin{prop} \label{prop:existsxepsfromsup}
Jika $S \subset \R$ tak kosong dan terbatas di atas,
maka untuk setiap $\epsilon > 0$ terdapat $x \in S$ sedemikian
sehingga $(\sup\, S) - \epsilon < x \leq \sup\, S$.
\end{prop}


Untuk semakin mempermudah penggunaan supremum dan infimum, kita mungkin ingin
menulis $\sup\, A$ dan $\inf\, A$ tanpa perlu mencemaskan apakah $A$
terbatas dan tak kosong.  Kita buat definisi-definisi alami berikut.

\begin{defn}
Misalkan $A \subset \R$ suatu himpunan.
\begin{enumerate}[(i)]
\item Jika $A$ kosong, maka $\sup\, A \coloneqq -\infty$.
\item Jika $A$ tidak terbatas di atas, maka $\sup\, A \coloneqq \infty$.
\item Jika $A$ kosong, maka $\inf\, A \coloneqq \infty$.
\item Jika $A$ tidak terbatas di bawah, maka $\inf\, A \coloneqq -\infty$.
\end{enumerate}
\end{defn}

Pernyataan-pernyataan dalam definisi di atas
muncul secara alami jika, alih-alih bekerja pada
himpunan terurut $\R$, kita mengambil supremum dan infimum
pada himpunan terurut
$\R^* \coloneqq \R \cup \{ -\infty , \infty\}$
\glsadd{not:extreal},\glsadd{not:infinity}
dengan menetapkan
\begin{equation*}
-\infty < \infty \quad \text{dan} \quad
-\infty < x \quad \text{dan} \quad
x < \infty \quad \text{untuk setiap $x \in \R$}.
\end{equation*}
Himpunan $\R^*$ disebut himpunan \emph{\myindex{bilangan real diperluas}}.
Untuk kemudahan, $\infty$ dan $-\infty$ terkadang diperlakukan seolah-olah merupakan
bilangan, tetapi kita tidak mengizinkan operasi aritmetika sebarang terhadapnya.
Kita dapat mendefinisikan beberapa operasi aritmetika pada $\R^*$.  Sebagian besar operasi
diperluas dengan cara yang jelas, tetapi kita harus membiarkan
$\infty-\infty$, $0 \cdot (\pm\infty)$, dan $\frac{\pm\infty}{\pm\infty}$
tidak terdefinisi.
Kita tidak
menggunakan aritmetika ini;
aritmetika tersebut mudah menimbulkan kekeliruan karena $\R^*$ bukan lapangan.
Sekarang kita dapat mengambil supremum dan infimum tanpa khawatir akan kekosongan atau
ketakterbatasan.  Dalam buku ini, kita umumnya menghindari
penggunaan $\R^*$ di luar latihan dan menyerahkan generalisasi semacam itu kepada pembaca yang berminat.

\subsection{Maksimum dan minimum}

Berdasarkan \exerciseref{exercise:finitesethasminmax},
suatu himpunan bilangan yang hingga selalu memiliki supremum dan infimum,
dan keduanya termuat dalam himpunan itu sendiri.
Dalam kasus ini, biasanya kita tidak menggunakan istilah supremum atau infimum.
Jika suatu himpunan bilangan real $A$ terbatas di atas dan
$\sup\, A \in A$, kita dapat menggunakan istilah
\emph{\myindex{maksimum}} dan notasi $\max\, A$\glsadd{not:max} untuk menyatakan supremum tersebut.
Demikian pula untuk infimum: Jika $A$ terbatas di bawah
dan $\inf\, A \in A$, kita dapat menggunakan
istilah \emph{\myindex{minimum}} dan notasi $\min\, A$.\glsadd{not:min}
Sebagai contoh,
\begin{align*}
& \max \{ 1,2.4,\pi,100 \} = 100 , \\
& \min \{ 1,2.4,\pi,100 \} = 1 .
\end{align*}
Walaupun menulis $\sup$ dan $\inf$ secara teknis mungkin
benar dalam situasi ini, $\max$ dan
$\min$ umumnya digunakan untuk menekankan bahwa supremum atau infimum
berada dalam himpunan itu sendiri, khususnya ketika himpunannya hingga.

\subsection{Latihan}

\begin{exercise}
Buktikan bahwa
jika $t > 0$ ($t \in \R$), maka terdapat $n \in \N$ sedemikian sehingga
$\dfrac{1}{n^2} < t$.
\end{exercise}

\begin{exercise}
Buktikan bahwa
jika $t \geq 0$ ($t \in \R$), maka terdapat $n \in \N$ sedemikian sehingga $n-1 \leq t < n$.
\end{exercise}

\begin{exercise}
Selesaikan bukti \propref{prop:supinfalg}.
\end{exercise}

\begin{exercise}
Misalkan $x, y \in \R$.  Andaikan $x^2 + y^2 = 0$.  Buktikan bahwa $x = 0$ dan $y = 0$.
\end{exercise}

\begin{exercise}
Tunjukkan bahwa $\sqrt{3}$ irasional.
\end{exercise}

\begin{exercise}
Misalkan $n \in \N$.
Tunjukkan bahwa $\sqrt{n}$ merupakan bilangan bulat atau
irasional.
\end{exercise}

\begin{exercise}
Buktikan \emph{\myindex{ketaksamaan rata-rata aritmetika-geometrik}}.
Untuk dua bilangan real positif $x,y$,
\begin{equation*}
\sqrt{xy} \leq \frac{x+y}{2} .
\end{equation*}
Lebih lanjut, kesamaan berlaku jika dan hanya jika $x=y$.
\end{exercise}

\begin{exercise}
Tunjukkan bahwa untuk setiap pasangan bilangan real $x$ dan $y$ sedemikian sehingga $x < y$,
terdapat bilangan irasional $s$ sedemikian sehingga $x < s < y$.  Petunjuk:
Terapkan kerapatan $\Q$ pada $\dfrac{x}{\sqrt{2}}$ dan
$\dfrac{y}{\sqrt{2}}$, tetapi hindari nol.
\end{exercise}

\begin{exercise} \label{exercise:supofsum}
Misalkan $A$ dan $B$ merupakan dua himpunan bilangan real tak kosong dan terbatas.  Misalkan
$C \coloneqq \{ a+b : a \in A, b \in B \}$.
Tunjukkan bahwa $C$ merupakan himpunan terbatas dan bahwa
\begin{equation*}
\sup\,C = \sup\,A + \sup\,B
\qquad \text{dan} \qquad
\inf\,C = \inf\,A + \inf\,B .
\end{equation*}
\end{exercise}

\begin{exercise}
Misalkan $A$ dan $B$ merupakan dua himpunan bilangan real nonnegatif yang tak kosong dan terbatas.
Definisikan himpunan
$C \coloneqq \{ ab : a \in A, b \in B \}$.
Tunjukkan bahwa $C$ merupakan himpunan terbatas dan bahwa
\begin{equation*}
\sup\,C = (\sup\,A )( \sup\,B)
\qquad \text{dan} \qquad
\inf\,C = (\inf\,A )( \inf\,B).
\end{equation*}
\end{exercise}

\begin{exercise}[Sulit] \label{exercise:rootexistshard}
Diberikan $x > 0$ dan $n \in \N$, tunjukkan bahwa terdapat bilangan real positif unik
$r$ sedemikian sehingga $x = r^n$.  Biasanya, $r$ dilambangkan dengan $x^{1/n}$.
\end{exercise}

\begin{exercise}[Mudah]
Buktikan \propref{prop:existsxepsfromsup}.
\end{exercise}

\begin{exercise} \label{exercise:bernoulliineq}
Buktikan \emph{\myindex{ketaksamaan Bernoulli}}%
\footnote{%
Dinamai menurut matematikawan Swiss
\href{https://en.wikipedia.org/wiki/Jacob_Bernoulli}{Jacob Bernoulli}
(1655--1705).}%
: Jika $1+x > 0$, maka
untuk setiap $n \in \N$, berlaku $(1+x)^n \geq 1+nx$.
\end{exercise}

\begin{exercise} \label{exercise:sqrt2QorR}
Buktikan $\sup \{ x \in \Q : x^2 < 2 \} = \sup \{ x \in \R : x^2 < 2 \}$.
\end{exercise}

\begin{exercise} \label{exercise:Dedekind}
\leavevmode
\begin{enumerate}[a)]
\item
Buktikan bahwa untuk setiap $y \in \R$, berlaku $\sup \{ x \in \Q : x < y \} = y$.
\item
Misalkan $A \subset \Q$ suatu himpunan tak kosong yang terbatas di atas sedemikian sehingga setiap kali $x
\in A$ dan $t \in \Q$ dengan $t < x$, maka $t \in A$.  Andaikan pula
$\sup\, A \notin A$.  Tunjukkan bahwa terdapat $y \in \R$ sedemikian sehingga
$A = \{ x \in \Q : x < y \}$.  Himpunan seperti $A$ disebut
\emph{\myindex{potongan Dedekind}}.
\item
Tunjukkan bahwa terdapat bijeksi antara $\R$ dan potongan-potongan Dedekind.
\end{enumerate}
Catatan: Dedekind menggunakan himpunan seperti dalam bagian b) pada konstruksi
bilangan real yang dikembangkannya.
\end{exercise}

\begin{exercise}
Buktikan bahwa jika $A \subset \Z$ merupakan himpunan bagian tak kosong yang terbatas di bawah, maka
terdapat unsur terkecil dalam $A$.
Kemudian, jelaskan mengapa pernyataan ini akan menyederhanakan bukti
\thmref{thm:arch} bagian \ref{thm:arch:ii} sehingga kita tidak perlu mengandaikan
$x \geq 0$.
\end{exercise}

\begin{samepage}
\begin{exercise} \label{exercise:realpower}
Andaikan kita mengetahui bahwa $x^{1/n}$ ada untuk setiap $x > 0$ dan setiap $n \in
\N$ (lihat \exerciseref{exercise:rootexistshard} di atas).
Untuk bilangan bulat $p$ dan $q > 0$ dengan $\nicefrac{p}{q}$ dalam bentuk paling sederhana,
definisikan $x^{p/q} \coloneqq
{(x^{1/q})}^p$.
\begin{enumerate}[a)]
\item
Tunjukkan bahwa pangkat tersebut terdefinisi dengan baik sekalipun
pecahannya tidak dalam bentuk paling sederhana: Jika $\nicefrac{p}{q} =
\nicefrac{m}{k}$ dengan $m$ dan $k > 0$ bilangan bulat, maka
${(x^{1/q})}^p = {(x^{1/k})}^m$.
\item
Misalkan $x$ dan $y$ dua bilangan positif dan $r$ suatu bilangan rasional.
Dengan mengandaikan $r > 0$, tunjukkan bahwa
$x < y$ jika dan hanya jika $x^r < y^r$.  Kemudian, andaikan $r < 0$ dan tunjukkan:
$x < y$ jika dan hanya jika $x^r > y^r$.
\item
Andaikan $x > 1$ dan $r,s$ bilangan-bilangan rasional dengan $r < s$.
Tunjukkan bahwa $x^r < x^s$.  Jika $0 < x < 1$ dan $r < s$, tunjukkan bahwa $x^r > x^s$.
Petunjuk: Tuliskan $r$ dan $s$ dengan penyebut yang sama.
\item
(Menantang)%
\footnote{Dalam
\sectionref{sec:logandexp}
kita akan mendefinisikan fungsi eksponensial dan logaritma, serta
mendefinisikan $x^z \coloneqq \exp(z \ln x)$.  Dengan demikian, kita akan memiliki
perangkat yang memadai untuk membuat pembuktian pernyataan-pernyataan ini jauh lebih mudah.  Namun, pada
tahap ini kita belum memiliki perangkat tersebut.}
Untuk bilangan irasional $z \in \R \setminus \Q$ dan $x > 1$, definisikan
$x^z \coloneqq \sup \{ x^r : r \leq z, r \in \Q \}$,
untuk $x=1$, definisikan $1^z = 1$,
dan untuk $0 < x < 1$, definisikan
$x^z \coloneqq \inf \{ x^r : r \leq z, r \in \Q \}$.
Buktikan kedua pernyataan pada bagian b) untuk setiap bilangan real $z$.
\end{enumerate}
\end{exercise}
\end{samepage}



%%%%%%%%%%%%%%%%%%%%%%%%%%%%%%%%%%%%%%%%%%%%%%%%%%%%%%%%%%%%%%%%%%%%%%%%%%%%%%

\sectionnewpage
\section{Nilai mutlak dan fungsi terbatas} \label{sec:absval}

\sectionnotes{0,5--1 pertemuan kuliah}

Salah satu konsep yang akan kita jumpai berulang kali adalah
\emph{\myindex{nilai mutlak}}.
Nilai mutlak dapat dipandang sebagai \myquote{ukuran} suatu bilangan real.
Berikut definisi formalnya:
\glsadd{not:abs}
\begin{equation*}
\sabs{x} \coloneqq
\begin{cases}
x & \text{jika } x \geq 0, \\
-x & \text{jika } x < 0 .
\end{cases}
\end{equation*}
Proposisi berikut memberikan
sifat-sifat utama nilai
mutlak.

\begin{prop} \label{prop:absbas}
\leavevmode
\begin{enumerate}[(i)]
\item \label{prop:absbas:i} $\sabs{x} \geq 0$; lebih lanjut, $\sabs{x}=0$ jika dan hanya jika $x = 0$.
\item \label{prop:absbas:ii} $\sabs{-x} = \sabs{x}$ untuk setiap $x \in \R$.
\item \label{prop:absbas:iii} $\sabs{xy} = \sabs{x}\sabs{y}$ untuk setiap $x,y \in \R$.
\item \label{prop:absbas:iv} $\sabs{x}^2 = x^2$ untuk setiap $x \in \R$.
\item \label{prop:absbas:v} $\sabs{x} \leq y$ jika dan hanya jika $-y \leq x \leq y$.
\item \label{prop:absbas:vi} $-\sabs{x} \leq x \leq \sabs{x}$ untuk setiap $x \in \R$.
\end{enumerate}
\end{prop}

\begin{proof}
\ref{prop:absbas:i}:
Pertama, andaikan $x \geq 0$.
Maka $\sabs{x} = x \geq 0$.  Selain itu, $\sabs{x} = x = 0$ jika dan hanya
jika $x=0$.
Di sisi lain, jika $x < 0$, maka $\sabs{x} = -x > 0$, dan $\sabs{x}$ tidak pernah nol.

\medskip

\ref{prop:absbas:ii}: Jika $x > 0$, maka $-x < 0$ sehingga $\sabs{-x} = -(-x) = x =
\sabs{x}$.  Demikian pula ketika $x < 0$, atau $x = 0$.

\medskip

\ref{prop:absbas:iii}:
Jika $x$ atau $y$ nol, hasilnya langsung jelas.  Ketika $x$ dan
$y$ keduanya positif, maka $\sabs{x}\sabs{y} = xy$.  Hasil kali bilangan-bilangan
positif bernilai positif, sehingga $xy$ bernilai
positif dan $xy = \sabs{xy}$.
Jika $x$ dan $y$ keduanya negatif,
maka $xy=(-x)(-y)$ bernilai positif karena alasan yang sama dan sekali lagi $\sabs{xy}=xy$.  Selain itu,
$\sabs{x}\sabs{y} = (-x)(-y) = xy$.
Jadi, hasilnya berlaku jika $x$ dan $y$ bertanda sama.
Jika $x > 0$ dan $y < 0$, maka $\sabs{x}\sabs{y} = x(-y) = -(xy)$.  Sekarang
$xy$ bernilai negatif dan $\sabs{xy} = -(xy)$, sehingga hasilnya berlaku.  Demikian pula ketika
$x < 0$ dan $y > 0$.

\medskip

\ref{prop:absbas:iv}:
Langsung jelas jika $x \geq 0$.  Jika $x < 0$, maka $\sabs{x}^2 = {(-x)}^2 =
x^2$.

\medskip

\ref{prop:absbas:v}:
Andaikan $\sabs{x} \leq y$.
Jika $x \geq 0$, maka $x \leq y$.
Akibatnya, $y \geq 0$, sehingga $-y \leq 0 \leq x$.
Jadi, $-y \leq x \leq y$ berlaku.
Jika $x < 0$, maka $\sabs{x} \leq y$ berarti $-x \leq y$.
Dengan mengalikan kedua ruas dengan $-1$, kita memperoleh $x \geq -y$.
Sekali lagi $y \geq 0$, sehingga $y \geq 0 > x$.
Dengan demikian, $-y \leq x \leq y$.

Di sisi lain, andaikan $-y \leq x \leq y$ berlaku.
Jika $x \geq 0$, maka $x \leq y$ ekuivalen dengan $\sabs{x} \leq y$.
Jika $x < 0$, maka $-y \leq x$ mengakibatkan $(-x) \leq y$,
yang ekuivalen dengan $\sabs{x} \leq y$.

\medskip

\ref{prop:absbas:vi}:
Terapkan \ref{prop:absbas:v} dengan $y = \sabs{x}$.
\end{proof}

Salah satu sifat yang cukup sering digunakan sehingga diberi nama adalah
\emph{\myindex{ketaksamaan segitiga}}.

\begin{prop}[Ketaksamaan Segitiga]
$\sabs{x+y} \leq \sabs{x}+\sabs{y}$
untuk setiap $x, y \in \R$.
\end{prop}

\begin{proof}
\propref{prop:absbas} memberikan
$- \sabs{x} \leq x \leq \sabs{x}$ dan
$- \sabs{y} \leq y \leq \sabs{y}$.  Jumlahkan kedua ketaksamaan ini untuk memperoleh
\begin{equation*}
- \bigl(\sabs{x}+\sabs{y}\bigr) \leq x+y \leq \sabs{x}+ \sabs{y} .
\end{equation*}
Terapkan kembali \propref{prop:absbas} untuk memperoleh
$\sabs{x+y} \leq \sabs{x}+\sabs{y}$.
\end{proof}

Terdapat versi-versi lain dari ketaksamaan segitiga yang juga sering digunakan.

\begin{cor}
Misalkan $x,y \in \R$.
\begin{enumerate}[(i)]
\item \emph{(\myindex{ketaksamaan segitiga terbalik})}
~
$\babs{ \, \sabs{x}-\sabs{y} \, } \leq \sabs{x-y}$.
\item $\sabs{x-y} \leq \sabs{x}+\sabs{y}$.
\end{enumerate}
\end{cor}

\begin{proof}
Masukkan $x=a-b$ dan $y=b$ ke dalam
ketaksamaan segitiga baku untuk memperoleh
\begin{equation*}
\sabs{a} = \sabs{a-b+b} \leq \sabs{a-b} + \sabs{b} ,
\end{equation*}
atau $\sabs{a}-\sabs{b} \leq \sabs{a-b}$.  Dengan menukar peran $a$ dan $b$,
kita memperoleh
$\sabs{b}-\sabs{a} \leq \sabs{b-a} = \sabs{a-b}$.  Dengan menerapkan
\propref{prop:absbas}, kita memperoleh ketaksamaan segitiga
terbalik.

Butir kedua dalam akibat tersebut diperoleh dari
ketaksamaan segitiga baku dengan mengganti $y$ dengan $-y$, dan memperhatikan bahwa $\sabs{-y} =
\sabs{y}$.
\end{proof}

\begin{cor}
Misalkan $x_1, x_2, \ldots, x_n \in \R$.  Maka
\begin{equation*}
\sabs{x_1 + x_2 + \cdots + x_n} \leq
\sabs{x_1} + \sabs{x_2} + \cdots + \sabs{x_n} .
\end{equation*}
\end{cor}

\begin{proof}
Kita menggunakan \hyperref[induction:thm]{induksi}.
Kesimpulannya berlaku secara langsung untuk $n=1$, dan
untuk $n=2$ merupakan ketaksamaan segitiga baku.  Andaikan akibat tersebut
berlaku untuk $n$.  Pertimbangkan $n+1$ bilangan $x_1,x_2,\ldots,x_{n+1}$:
Pertama, gunakan ketaksamaan segitiga baku, kemudian gunakan hipotesis
induksi
\begin{equation*}
\begin{split}
\sabs{x_1 + x_2 + \cdots + x_n + x_{n+1}} & \leq
\sabs{x_1 + x_2 + \cdots + x_n} + \sabs{x_{n+1}} \\
& \leq
\sabs{x_1} + \sabs{x_2} + \cdots + \sabs{x_n} + \sabs{x_{n+1}} .  \qedhere
\end{split}
\end{equation*}
\end{proof}

Mari kita lihat contoh penggunaan ketaksamaan segitiga.

\begin{example}
Temukan bilangan $M$ sedemikian sehingga $\sabs{x^2-9x+1} \leq M$ untuk setiap $-1 \leq x \leq
5$.

Dengan menggunakan ketaksamaan segitiga, tuliskan
\begin{equation*}
\sabs{x^2-9x+1} \leq \sabs{x^2}+\sabs{-9x}+\sabs{1}
=
\sabs{x}^2+9\sabs{x}+1 .
\end{equation*}
Ekspresi
$\sabs{x}^2+9\sabs{x}+1$ bernilai paling besar ketika $\sabs{x}$ paling besar (mengapa?).  Pada interval
yang diberikan, $\sabs{x}$ paling besar ketika $x=5$, sehingga $\sabs{x}=5$.  Salah satu
nilai yang mungkin untuk $M$ adalah
\begin{equation*}
M = 5^2+9(5)+1 = 71 .
\end{equation*}
Tentu saja, terdapat nilai-nilai $M$ lain yang memenuhi.  Batas 71
jauh lebih besar daripada yang
diperlukan, tetapi kita tidak meminta nilai $M$ terbaik, hanya satu nilai yang memenuhi.
\end{example}

Contoh terakhir mengantarkan kita pada konsep fungsi terbatas.

\begin{defn}
Andaikan $f \colon D \to \R$ suatu fungsi.  Kita mengatakan bahwa $f$
\emph{terbatas}\index{fungsi terbatas}
jika terdapat bilangan $M$
sedemikian sehingga $\babs{f(x)} \leq M$ untuk setiap $x \in D$.
\end{defn}

Dalam contoh tersebut, kita membuktikan bahwa $x^2-9x+1$ terbatas ketika dipandang sebagai
fungsi pada $D = \{ x : -1 \leq x \leq 5 \}$.  Di sisi lain,
jika kita memandang polinom yang sama sebagai fungsi pada seluruh garis bilangan real $\R$,
maka fungsi tersebut tidak terbatas.

\begin{myfigureht}
\myincludegraphics{boundedfunc}{%
Diagram grafik suatu fungsi.
Domain D ditandai pada sumbu horizontal, dan digambar sebuah garis sinusoidal
yang merepresentasikan grafik fungsi pada domain tersebut.  Nilai maksimum
fungsi digambar sebagai garis putus-putus horizontal yang menyentuh
grafik fungsi dan diberi label sup f dari D.  Demikian pula, nilai
minimum diberi label inf f dari D.  Dua garis putus-titik lainnya diberi label
M dan minus M.  Garis berlabel M berada di atas grafik dan garis berlabel
minus M berada di bawah grafik.  Pada sumbu vertikal, himpunan nilai
yang dicapai fungsi diberi label f dari D.}
\caption{Contoh fungsi terbatas, suatu batas $M$, serta supremum dan infimumnya.\label{boundedfuncfig}}
\end{myfigureht}
Untuk fungsi $f \colon D \to \R$, kita menulis
(lihat \figureref{boundedfuncfig} untuk
sebuah contoh)
\glsadd{not:supf}\glsadd{not:inff}
\begin{equation*}
\sup_{x \in D} f(x) \coloneqq \sup\, f(D) \qquad \text{dan} \qquad
\inf_{x \in D} f(x) \coloneqq \inf\, f(D) .
\end{equation*}
Terkadang kita juga mengganti \myquote{$x \in D$} dengan suatu ekspresi.
Sebagai contoh, jika, seperti sebelumnya, $f(x) = x^2-9x+1$, untuk $-1 \leq x \leq 5$,
sedikit kalkulus menunjukkan bahwa
\begin{equation*}
\sup_{x \in D} f(x) =
\sup_{-1 \leq x \leq 5} ( x^2 -9x+1 ) = 11,
\qquad
\inf_{x \in D} f(x) =
\inf_{-1 \leq x \leq 5} ( x^2 -9x+1 ) = \nicefrac{-77}{4} .
\end{equation*}



\begin{prop} \label{prop:funcsupinf}
Jika $f \colon D \to \R$ dan $g \colon D \to \R$ ($D$ tak kosong) merupakan
fungsi-fungsi terbatas\footnote{Hipotesis keterbatasan digunakan demi kesederhanaan;
hipotesis ini dapat dihapus jika kita menggunakan bilangan real diperluas.}
dan
\begin{equation*}
f(x) \leq g(x) \qquad \text{untuk setiap } x \in D,
\end{equation*}
maka
\begin{equation} \label{prop:funcsupinf:eq}
\sup_{x \in D} f(x) \leq \sup_{x \in D} g(x)
\qquad \text{dan} \qquad
\inf_{x \in D} f(x) \leq \inf_{x \in D} g(x) .
\end{equation}
\end{prop}

Berhati-hatilah dengan variabel-variabelnya.  Variabel $x$ pada ruas kiri
ketaksamaan dalam \eqref{prop:funcsupinf:eq}
berbeda dari variabel $x$ pada ruas kanan.  Sebagai contoh,
ketaksamaan pertama sebaiknya dipandang sebagai
\begin{equation*}
\sup_{x \in D} f(x) \leq \sup_{y \in D} g(y) .
\end{equation*}
Mari kita buktikan ketaksamaan ini.  Jika $b$ merupakan batas atas bagi $g(D)$, maka
$f(x) \leq g(x) \leq b$ untuk setiap $x \in D$, sehingga $b$ juga merupakan batas
atas bagi $f(D)$,
atau $f(x) \leq b$ untuk setiap $x \in D$.
Ambil batas atas terkecil dari $g(D)$; maka untuk setiap $x \in D$ diperoleh
\begin{equation*}
f(x) \leq \sup_{y \in D} g(y) .
\end{equation*}
Oleh karena itu,
$\sup_{y \in D} g(y)$ merupakan batas atas bagi $f(D)$, sehingga lebih besar atau
sama dengan batas atas terkecil dari $f(D)$.
\begin{equation*}
\sup_{x \in D} f(x) \leq \sup_{y \in D} g(y) .
\end{equation*}
Ketaksamaan kedua (pernyataan tentang infimum) diserahkan sebagai latihan
(\exerciseref{exercise:finishfuncsupinf}).

\medskip

Kekeliruan yang umum adalah menyimpulkan
\begin{equation} \label{rn:av:ltnottrue}
\sup_{x \in D} f(x) \leq \inf_{y \in D} g(y) .
\end{equation}
Ketaksamaan \eqref{rn:av:ltnottrue} tidak benar berdasarkan hipotesis
proposisi di atas.  Untuk ketaksamaan yang lebih kuat ini,
kita memerlukan hipotesis yang lebih kuat
\begin{equation*}
f(x) \leq g(y) \qquad \text{untuk setiap } x \in D \text{ dan } y \in D.
\end{equation*}
Bukti serta sebuah contoh penyangkal diserahkan sebagai latihan
(\exerciseref{exercise:supinffuncineq}).

\subsection{Latihan}

\begin{exercise}
Tunjukkan bahwa
$\sabs{x-y} < \epsilon$ jika dan hanya jika $x-\epsilon < y < x+\epsilon$.
\end{exercise}

\begin{exercise}
Tunjukkan: \qquad
%\begin{enumerate}[a)]
%\item
a)
$\max \{x,y\} = \dfrac{x+y+\sabs{x-y}}{2}$
%\item
\qquad
b)
$\min \{x,y\} = \dfrac{x+y-\sabs{x-y}}{2}$
%\end{enumerate}
\end{exercise}

\begin{exercise}
Temukan bilangan $M$ sedemikian sehingga $\sabs{x^3-x^2+8x} \leq M$ untuk setiap $-2 \leq x \leq
10$.
\end{exercise}

\begin{exercise} \label{exercise:finishfuncsupinf}
Selesaikan bukti \propref{prop:funcsupinf}.
Artinya,
diberikan suatu himpunan $D$,
dan dua fungsi terbatas
$f \colon D \to \R$ dan $g \colon D \to \R$
sedemikian sehingga $f(x) \leq g(x)$ untuk setiap $x \in D$, buktikan bahwa
\begin{equation*}
\inf_{x\in D} f(x) \leq \inf_{x\in D} g(x) .
\end{equation*}
\end{exercise}

\begin{exercise} \label{exercise:supinffuncineq}
Misalkan
$f \colon D \to \R$ dan $g \colon D \to \R$ merupakan fungsi-fungsi ($D$ tak kosong).
\begin{enumerate}[a)]
%
\item
Andaikan
$f(x) \leq g(y)$ untuk setiap $x \in D$ dan $y \in D$.  Tunjukkan bahwa
\begin{equation*}
\sup_{x\in D} f(x) \leq \inf_{x\in D} g(x) .
\end{equation*}
%
\item
Temukan $D$, $f$, dan $g$ tertentu sedemikian sehingga
$f(x) \leq g(x)$ untuk setiap $x \in D$, tetapi
\begin{equation*}
\sup_{x\in D} f(x) > \inf_{x\in D} g(x) .
\end{equation*}
\end{enumerate}
\end{exercise}

\begin{exercise}
Buktikan \propref{prop:funcsupinf} tanpa asumsi bahwa
fungsi-fungsinya terbatas.  Petunjuk: Anda perlu menggunakan bilangan
real diperluas.
\end{exercise}

\begin{exercise} \label{exercise:sumofsup}
\pagebreak[2]
Misalkan $D$ suatu himpunan tak kosong.
Andaikan $f \colon D \to \R$ dan $g \colon D \to \R$ merupakan fungsi-fungsi terbatas.
\begin{enumerate}[a)]
\item
Tunjukkan
\begin{equation*}
%mbxlatex \begin{aligned}
%mbxlatex &
\sup_{x\in D} \bigl(f(x) + g(x) \bigr) \leq
\sup_{x\in D} f(x)
+
\sup_{x\in D} g(x)
\qquad \text{dan}
%mbxSTARTIGNORE
\qquad
%mbxENDIGNORE
%mbxlatex \\
%mbxlatex &
\inf_{x\in D} \bigl(f(x) + g(x) \bigr) \geq
\inf_{x\in D} f(x)
+
\inf_{x\in D} g(x) .
%mbxlatex \end{aligned}
\end{equation*}
\item
Temukan sebuah contoh (atau beberapa contoh) yang menghasilkan ketaksamaan ketat.
\end{enumerate}
\end{exercise}

\begin{exercise}
\pagebreak[2]
Andaikan $D$ tak kosong, $f \colon D \to \R$ dan $g \colon D \to \R$ merupakan
fungsi-fungsi terbatas, dan
$\alpha \in \R$.
\begin{enumerate}[a)]
\item
Tunjukkan bahwa $\alpha f \colon D \to \R$ yang didefinisikan oleh $(\alpha f) (x) \coloneqq \alpha
f(x)$ merupakan fungsi terbatas.
\item
Tunjukkan bahwa $f+g \colon D \to \R$ yang didefinisikan oleh $(f+g) (x) \coloneqq f(x) + g(x)$
merupakan fungsi terbatas.
\end{enumerate}
\end{exercise}

\begin{exercise}
Misalkan $f \colon D \to \R$ dan $g \colon D \to \R$ merupakan fungsi-fungsi dengan $D$
tak kosong, $\alpha \in \R$, dan
ingat kembali arti $f+g$ dan $\alpha f$ dari latihan sebelumnya.
\begin{enumerate}[a)]
\item
Buktikan bahwa jika $f+g$ dan $g$ terbatas, maka $f$ terbatas.
\item
Temukan contoh dengan $f$ dan $g$ keduanya tak terbatas, tetapi $f+g$
terbatas.
\item
Buktikan bahwa jika $f$ terbatas tetapi $g$ tak terbatas, maka $f+g$
tak terbatas.
\item
Temukan contoh dengan $f$ tak terbatas tetapi $\alpha f$ terbatas.
\end{enumerate}
\end{exercise}

%%%%%%%%%%%%%%%%%%%%%%%%%%%%%%%%%%%%%%%%%%%%%%%%%%%%%%%%%%%%%%%%%%%%%%%%%%%%%%

\sectionnewpage
\section{Interval dan ukuran \texorpdfstring{$\R$}{R}}
\label{sec:intandsizeR}

\sectionnotes{0,5--1 pertemuan kuliah (bukti bahwa $\R$ tak terhitung dapat dijadikan opsional)}

Anda tentu pernah melihat notasi untuk interval\index{interval}
sebelumnya, tetapi mari kita berikan definisi
formal di sini.  Untuk $a, b \in \R$ sedemikian sehingga $a < b$, kita definisikan
\glsadd{not:openbdinterval}%
\glsadd{not:closedbdinterval}%
\glsadd{not:halfopenbdinterval}%
\glsadd{not:openubdinterval}%
\glsadd{not:closedubdinterval}%
\begin{align*}
& [a,b] \coloneqq \{ x \in \R : a \leq x \leq b \}, \\
& (a,b) \coloneqq \{ x \in \R : a < x < b \}, \\
& (a,b] \coloneqq \{ x \in \R : a < x \leq b \}, \\
& [a,b) \coloneqq \{ x \in \R : a \leq x < b \} .
\end{align*}
Interval $[a,b]$ disebut \emph{\myindex{interval tertutup}}
dan $(a,b)$ disebut \emph{\myindex{interval terbuka}}.  Interval-interval
berbentuk $(a,b]$ dan $[a,b)$ disebut
\emph{interval setengah terbuka}\index{interval setengah terbuka}.

Interval-interval di atas adalah \emph{interval terbatas}\index{interval
terbatas}, karena $a$ dan $b$ keduanya bilangan real.\footnote{Ketika kita menulis
$(a,b)$, $[a,b]$, dan seterusnya, yang kita maksud selalu interval terbatas kecuali jika secara eksplisit dikatakan
bahwa $a$ atau $b$ dapat berupa $-\infty$ atau $\infty$.}
Kita mendefinisikan \emph{interval tak terbatas}\index{interval tak terbatas},
\begin{align*}
& [a,\infty) \coloneqq \{ x \in \R : a \leq x \}, \\
& (a,\infty) \coloneqq \{ x \in \R : a < x \}, \\
& (-\infty,b] \coloneqq \{ x \in \R : x \leq b \}, \\
& (-\infty,b) \coloneqq \{ x \in \R : x < b \} .
\end{align*}
Demi kelengkapan, kita definisikan $(-\infty,\infty) \coloneqq \R$.
Interval $[a,\infty)$, $(-\infty,b]$, dan $\R$ terkadang disebut
\emph{\myindex{interval tertutup tak terbatas}},
sedangkan $(a,\infty)$, $(-\infty,b)$, dan $\R$ terkadang disebut
\emph{\myindex{interval terbuka tak terbatas}}.

Bukti proposisi berikut diserahkan sebagai latihan.
Singkatnya, interval adalah himpunan dengan setidaknya dua titik
yang memuat semua titik di antara sebarang dua
titiknya.\footnote{Terkadang himpunan satu titik dan himpunan kosong
juga disebut interval, tetapi dalam buku ini interval memiliki setidaknya dua titik.
Artinya, kita hanya mendefinisikan interval terbatas jika $a < b$.}

\begin{prop} \label{prop:intervaldef}
Suatu himpunan $I \subset \R$ merupakan interval jika dan hanya jika
$I$ memuat setidaknya dua titik dan
untuk setiap $a,c \in I$ dan $b \in \R$ sedemikian sehingga $a < b < c$, berlaku $b \in I$.
\end{prop}

Kita telah melihat bahwa setiap interval terbuka $(a,b)$ (tentu saja dengan $a < b$)
tak kosong.  Sebagai contoh, interval tersebut memuat bilangan $\frac{a+b}{2}$.
Fakta yang tidak terduga adalah bahwa dari sudut pandang teori himpunan,
semua interval memiliki \myquote{ukuran} yang sama, yakni semuanya memiliki
kardinalitas yang sama.  Sebagai contoh, pemetaan $f(x) \coloneqq 2x$
memetakan interval $[0,1]$ secara bijektif ke interval $[0,2]$.

Mungkin yang lebih menarik,
fungsi $f(x) \coloneqq \tan(x)$
merupakan pemetaan bijektif dari $(-\nicefrac{\pi}{2},\nicefrac{\pi}{2})$ ke $\R$.  Jadi, interval terbatas
$(-\nicefrac{\pi}{2},\nicefrac{\pi}{2})$ memiliki kardinalitas yang sama dengan $\R$.  Tidak mudah
mengonstruksi pemetaan bijektif dari $[0,1]$ ke
$(0,1)$, tetapi hal itu dapat dilakukan.

Dan jangan khawatir, memang ada cara untuk mengukur \myquote{ukuran} himpunan bagian
bilangan real yang
\myquote{mendeteksi}
perbedaan antara $[0,1]$ dan $[0,2]$.  Namun, definisi yang tepat
memerlukan perangkat yang jauh lebih banyak daripada yang kita miliki saat ini.

Mari kita bahas lebih lanjut kardinalitas interval dan, dengan demikian,
kardinalitas $\R$.  Kita telah melihat bahwa bilangan irasional ada;
artinya,
$\R \setminus \Q$ tak kosong.  Pertanyaannya: Ada berapa banyak bilangan irasional?
Ternyata bilangan irasional jauh lebih banyak daripada bilangan
rasional.  Kita telah melihat bahwa $\Q$ terhitung, dan kita akan menunjukkan
bahwa $\R$ tak terhitung.
Faktanya, kardinalitas $\R$ sama dengan
kardinalitas $\sP(\N)$, meskipun kita tidak akan membuktikan
klaim ini di sini.

\begin{thm}[Cantor]\index{teorema Cantor}
$\R$ tak terhitung.
\end{thm}

Kita menyajikan suatu versi
bukti asli Cantor dari tahun
1874 karena bukti ini memerlukan persiapan paling sedikit.  Biasanya bukti ini dinyatakan
sebagai kontradiksi.  Namun, bukti tersebut lebih mudah
dipahami sebagai bukti konstruktif langsung bahwa setiap himpunan bagian terhitung tak hingga
pasti tidak memuat setidaknya satu unsur.

\begin{proof}
Misalkan $X \subset \R$ suatu himpunan bagian
terhitung tak hingga sedemikian sehingga untuk setiap pasangan bilangan real
$a < b$, terdapat $x \in X$ sedemikian sehingga $a < x < b$.  Jika $\R$ terhitung,
kita dapat mengambil $X = \R$.  Kita akan menunjukkan bahwa $X$ pasti merupakan
himpunan bagian sejati, sehingga $X$ tidak mungkin sama dengan $\R$, dan $\R$ harus
tak terhitung.

Karena $X$ terhitung tak hingga,
terdapat bijeksi dari $\N$ ke $X$.  Kita menuliskan $X$ sebagai
barisan bilangan real $x_1, x_2, x_3,\ldots$, sedemikian sehingga
setiap bilangan dalam $X$
diberikan oleh $x_n$ untuk suatu $n \in \N$.

Secara induktif, kita
mengonstruksi dua barisan bilangan real $a_1,a_2,a_3,\ldots$ dan
$b_1,b_2,b_3,\ldots$.  Misalkan
$a_1 \coloneqq x_1$ dan $b_1 \coloneqq x_1+1$.
Perhatikan bahwa $a_1 < b_1$ dan
$x_1 \notin (a_1,b_1)$.
Untuk suatu $k > 1$,
andaikan $a_1,a_2,\ldots,a_{k-1}$ dan $b_1,b_2,\ldots,b_{k-1}$ telah
didefinisikan,
andaikan
$a_1 < a_2 < \cdots < a_{k-1} < b_{k-1} < \cdots < b_2 < b_1$,
dan andaikan
untuk setiap $j=1,2,\ldots,k-1$,
berlaku $x_\ell \notin (a_{j},b_{j})$
untuk $\ell=1,2,\ldots,j$.
\begin{enumerate}[(i)]
\item Definisikan $a_k \coloneqq x_n$, dengan $n$ dipilih sebagai bilangan terkecil di antara semua $n \in \N$
yang memenuhi $x_n \in (a_{k-1},b_{k-1})$.  Unsur $x_n$ semacam itu ada berdasarkan
asumsi kita mengenai $X$, dan $n \geq k$ berdasarkan asumsi mengenai $(a_{k-1},b_{k-1})$.
\item Selanjutnya, definisikan $b_k$ sebagai suatu bilangan real dalam $(a_{k},b_{k-1})$,
misalnya $b_k \coloneqq \frac{a_k+b_{k-1}}{2}$.
\end{enumerate}
Perhatikan bahwa %$a_k < b_k$ and
$a_{k-1} < a_k < b_k < b_{k-1}$.
Perhatikan pula bahwa $(a_{k},b_{k})$ tidak memuat $x_k$ sehingga
tidak memuat $x_{\ell}$ untuk $\ell=1,2,\ldots,k$.
Kedua barisan kini telah didefinisikan.

Klaim: \emph{$a_n < b_m$ untuk setiap $n$ dan $m$ dalam $\N$.} Bukti: Pertama,
andaikan $n < m$.  Maka $a_n < a_{n+1} < \cdots < a_{m-1} < a_m < b_m$.
Demikian pula untuk $n > m$.  Kasus dengan kedua indeks sama langsung mengikuti konstruksi.  Klaim terbukti.

Misalkan $A \coloneqq \{ a_n : n \in \N \}$ dan $B \coloneqq \{ b_n : n \in \N \}$.
Berdasarkan \propref{infsupineq:prop} dan klaim di atas,
\begin{equation*}
\sup\, A \leq \inf\, B .
\end{equation*}
Definisikan $y \coloneqq \sup\, A$.  Bilangan $y$ tidak mungkin merupakan unsur $A$: Jika $y = a_n$
untuk suatu $n$, maka $y < a_{n+1}$, yang mustahil.
Demikian pula, $y$ tidak mungkin merupakan unsur $B$.  Oleh karena itu,
$a_n < y$ untuk setiap $n\in \N$
dan $y < b_n$ untuk setiap $n\in \N$.
Dengan kata lain, untuk setiap $n \in \N$, berlaku $y \in (a_n,b_n)$.
Berdasarkan konstruksi barisan tersebut, $x_n \notin (a_n,b_n)$,
sehingga $y \neq x_n$.  Karena $y \neq x_n$ untuk setiap $n \in \N$, diperoleh $y
\notin X$.

Kita telah mengonstruksi bilangan real $y$ yang tidak berada di $X$, sehingga $X$
merupakan himpunan bagian sejati dari $\R$.
Barisan $x_1,x_2,\ldots$ tidak mungkin memuat semua unsur $\R$,
sehingga $\R$ tak terhitung.
\end{proof}

\subsection{Latihan}

\begin{exercise}
Untuk $a < b$, konstruksikan bijeksi eksplisit dari $(a,b]$ ke $(0,1]$.
\end{exercise}

\begin{exercise}
Misalkan $f \colon [0,1] \to (0,1)$ merupakan bijeksi.
Dengan menggunakan $f$, konstruksikan suatu
bijeksi dari $[-1,1]$ ke $\R$.
\end{exercise}

\begin{exercise} \label{exercise:intervaldef}
Buktikan \propref{prop:intervaldef}.
Artinya,
misalkan $I \subset \R$ merupakan himpunan bagian dengan setidaknya dua unsur
sedemikian sehingga jika $a < b < c$ dan $a, c \in I$, maka $b \in I$.
Buktikan bahwa $I$ merupakan salah satu dari sembilan jenis interval yang secara eksplisit
diberikan dalam bagian ini.
Selanjutnya, buktikan bahwa semua interval yang diberikan dalam bagian ini
memenuhi sifat tersebut.
\end{exercise}

\begin{exercise}[Sulit]
Konstruksikan bijeksi eksplisit dari $(0,1]$ ke $(0,1)$.
Petunjuk: Salah satu pendekatannya sebagai berikut: Pertama, petakan $(\nicefrac{1}{2},1]$
ke $(0,\nicefrac{1}{2}]$, lalu petakan
$(\nicefrac{1}{4},\nicefrac{1}{2}]$
ke $(\nicefrac{1}{2},\nicefrac{3}{4}]$, dan seterusnya.
Tuliskan pemetaan tersebut secara eksplisit, yaitu
tuliskan algoritme yang menentukan dengan tepat bilangan mana dipetakan ke mana.
Kemudian buktikan bahwa pemetaan tersebut merupakan bijeksi.
\end{exercise}

\begin{exercise}[Sulit]
Konstruksikan bijeksi eksplisit dari $[0,1]$ ke $(0,1)$.
\end{exercise}

\begin{exercise}
\leavevmode
\begin{enumerate}[a)]
\item
Tunjukkan bahwa setiap interval tertutup $[a,b]$ merupakan irisan terhitung
dari interval-interval terbuka.
\item
Tunjukkan bahwa setiap interval terbuka $(a,b)$
merupakan gabungan terhitung dari interval-interval tertutup.
\item
Tunjukkan bahwa irisan
dari suatu keluarga tak kosong dan mungkin tak hingga yang terdiri atas interval tertutup terbatas,
$\bigcap\limits_{\lambda \in I} [a_\lambda,b_\lambda]$,
adalah himpunan kosong, satu titik,
atau interval tertutup terbatas.
\end{enumerate}
\end{exercise}

\begin{exercise}
Misalkan $S$ merupakan himpunan interval terbuka yang saling lepas di $\R$.  Artinya,
jika $(a,b) \in S$ dan $(c,d) \in S$, maka $(a,b) = (c,d)$
atau $(a,b) \cap (c,d) = \emptyset$.  Buktikan bahwa $S$ merupakan himpunan terhitung.
\end{exercise}

\begin{exercise}
Buktikan bahwa kardinalitas $[0,1]$ sama dengan kardinalitas
$(0,1)$ dengan menunjukkan bahwa
$\babs{[0,1]} \leq \babs{(0,1)}$ dan
$\babs{(0,1)} \leq \babs{[0,1]}$.
Lihat \defnref{def:comparecards}.
Bukti ini memerlukan teorema Cantor--Bernstein--Schr\"oder, yang telah kita
nyatakan tanpa bukti.  Perhatikan bahwa bukti ini tidak menghasilkan
bijeksi eksplisit.
\end{exercise}

\begin{exercise}[Menantang]
Suatu bilangan $x$ disebut \emph{bilangan aljabar}\index{bilangan aljabar} jika $x$ merupakan akar dari
suatu polinomial tak nol dengan
koefisien bulat; dengan kata lain, $a_n x^n + a_{n-1} x^{n-1} + \cdots
+ a_1 x + a_0 = 0$ dengan $a_0,a_1,\ldots,a_n \in \Z$, tetapi tidak semuanya nol.
\begin{enumerate}[a)]
\item
Tunjukkan bahwa himpunan semua
bilangan aljabar terhitung.
\item
Tunjukkan bahwa terdapat bilangan yang bukan bilangan aljabar (bilangan transenden)
(ikuti jejak Cantor dan gunakan fakta bahwa $\R$ tak terhitung).
\end{enumerate}
Petunjuk: Anda boleh menggunakan fakta bahwa polinomial tak nol berderajat $n$ memiliki paling banyak $n$ akar
real.
\end{exercise}

\begin{exercise}[Menantang]
Misalkan $F$ merupakan himpunan semua fungsi $f \colon \R \to \R$.
Buktikan $\sabs{\R} < \sabs{F}$
dengan menggunakan hasil Cantor pada \thmref{cantorspowersetthm}.\footnote{Menariknya,
jika $C$ merupakan himpunan fungsi kontinu, maka $\sabs{\R} = \sabs{C}$.}
\end{exercise}

%%%%%%%%%%%%%%%%%%%%%%%%%%%%%%%%%%%%%%%%%%%%%%%%%%%%%%%%%%%%%%%%%%%%%%%%%%%%%%

\sectionnewpage
\section{Representasi desimal bilangan real}
\label{sec:decimals}

\sectionnotes{1 pertemuan kuliah (opsional)}

Kita sering membayangkan bilangan real melalui
\emph{representasi desimalnya}\index{representasi desimal}.
Yang dimaksud dengan
\emph{\myindex{digit}} (desimal)\index{digit desimal} adalah bilangan bulat
antara $0$ dan $9$.
Untuk bilangan bulat positif $n$, kita mencari digit-digit $d_K,d_{K-1},\ldots,d_2,d_1,d_0$ untuk suatu
$K$ (setiap $d_j$ merupakan bilangan bulat antara $0$ dan $9$) sedemikian sehingga
\begin{equation*}
n = d_K {10}^K + d_{K-1} {10}^{K-1} + \cdots + d_2 {10}^2 + d_1 10 + d_0 .
\end{equation*}
Kita sering mengasumsikan $d_K \neq 0$ (untuk menghindari nol di depan).
Untuk merepresentasikan $n$, kita menuliskan barisan
digit: $n = d_K d_{K-1} \cdots d_2 d_1 d_0$.

Dengan cara serupa, kita
merepresentasikan beberapa bilangan rasional. Untuk bilangan
$x$ tertentu, kita dapat mencari
bilangan bulat tak negatif $K$ dan $M$ serta digit-digit
$d_K,d_{K-1},\ldots,d_1,d_0,d_{-1},\ldots,d_{-M}$, sedemikian sehingga
\begin{equation*}
%mbxlatex \begin{aligned}
x
%mbxlatex &
= d_K {10}^K + d_{K-1} {10}^{K-1} + \cdots + d_2 {10}^2 + d_1 10 + d_0
%mbxlatex \\
%mbxlatex & \phantom{={}}
+ d_{-1} {10}^{-1} + d_{-2} {10}^{-2} + \cdots + d_{-M} {10}^{-M} .
%mbxlatex \end{aligned}
\end{equation*}
Kita menuliskan $x = d_K d_{K-1} \cdots d_1 d_0 \, . \, d_{-1} d_{-2} \cdots d_{-M}$.

Tidak setiap bilangan real memiliki representasi semacam itu.
Bahkan bilangan rasional sederhana
$\nicefrac{1}{3}$ tidak memilikinya. Bilangan irasional $\sqrt{2}$ juga
tidak memiliki representasi semacam itu. Untuk memperoleh representasi bagi
semua bilangan real, kita harus mengizinkan digit dalam jumlah tak hingga.

Tinjau hanya bilangan real dalam interval $[0,1]$. Jika
kita menemukan representasi untuk bilangan-bilangan ini, dengan menambahkan
bilangan bulat kepadanya kita memperoleh representasi untuk semua bilangan real.
Ambil suatu barisan digit desimal tak hingga:
\begin{equation*}
0.d_1d_2d_3\ldots.
\end{equation*}
Kita memiliki digit $d_j$ untuk setiap $j \in \N$.
Kita menomori ulang digit-digit tersebut untuk menghindari tanda negatif.
Kita tuliskan
\begin{equation*}
D_n \coloneqq
\frac{d_1}{10} +
\frac{d_2}{{10}^2} +
\frac{d_3}{{10}^3} +
\cdots +
\frac{d_n}{{10}^n} .
\end{equation*}
Kita katakan bahwa
barisan digit tersebut merepresentasikan bilangan real $x$ jika
\begin{equation*}
x =
\sup_{n \in \N} \left(
\frac{d_1}{10} +
\frac{d_2}{{10}^2} +
\frac{d_3}{{10}^3} +
\cdots +
\frac{d_n}{{10}^n}
\right) =
\sup_{n \in \N} \, D_n .
\end{equation*}
Bilangan $D_n=0.d_1 d_2 d_3 \ldots d_n$ adalah pemotongan $x$
hingga $n$ digit desimal.

\begin{prop} \label{prop:decimalprop}
\leavevmode
\begin{enumerate}[(i)]
\item
Setiap barisan digit tak hingga
$0.d_1d_2d_3\ldots$ merepresentasikan tepat satu bilangan real $x \in [0,1]$, dan
\begin{equation*}
D_n \leq x \leq D_n+\frac{1}{{10}^n} \qquad \text{untuk semua } n \in \N.
\end{equation*}
\item
Untuk setiap $x \in (0,1]$ terdapat barisan digit tak hingga
$0.d_1d_2d_3\ldots$ yang merepresentasikan $x$.
Terdapat representasi unik sedemikian sehingga
\begin{equation*}
D_n < x \leq D_n+\frac{1}{{10}^n} \qquad \text{untuk semua } n \in \N.
\end{equation*}
\end{enumerate}
\end{prop}

\begin{proof}
Kita mulai dengan butir pertama. Ambil sebarang barisan digit
tak hingga $0.d_1d_2d_3\ldots$. Dengan menggunakan rumus jumlah deret geometri, kita tuliskan
\begin{equation*}
\begin{split}
D_n =
\frac{d_1}{10} +
\frac{d_2}{{10}^2} +
\frac{d_3}{{10}^3} +
\cdots +
\frac{d_n}{{10}^n}
& \leq
\frac{9}{10} +
\frac{9}{{10}^2} +
\frac{9}{{10}^3} +
\cdots +
\frac{9}{{10}^n}
\\
& =
\frac{9}{10}
\bigl( 1 + \nicefrac{1}{10} +
{(\nicefrac{1}{10})}^2 + \cdots +
{(\nicefrac{1}{10})}^{n-1} \bigr)
\\
& =
\frac{9}{10}
\left(
\frac{1-{(\nicefrac{1}{10})}^{n}}{1-\nicefrac{1}{10}}
\right)
= 1-{(\nicefrac{1}{10})}^{n}
< 1 .
\end{split}
\end{equation*}
Secara khusus, $D_n < 1$ untuk semua $n$. Jumlah bilangan-bilangan tak negatif
bersifat tak negatif, sehingga $D_n \geq 0$, dan akibatnya
\begin{equation*}
0 \leq \sup_{n\in \N} \, D_n \leq 1 .
\end{equation*}
Oleh karena itu, $0.d_1d_2d_3\ldots$ merepresentasikan tepat satu bilangan real $x \coloneqq \sup_{n\in
\N} D_n \in [0,1]$.
Karena $x$ merupakan supremum, $D_n \leq x$.
Ambil $m \in \N$. Jika $m < n$, maka $D_m - D_n \leq 0$. Jika kedua indeks sama, selisih tersebut nol. Jika $m > n$, dengan
perhitungan seperti di atas diperoleh
\begin{equation*}
D_m - D_n
=
\frac{d_{n+1}}{{10}^{n+1}} +
\frac{d_{n+2}}{{10}^{n+2}} +
\frac{d_{n+3}}{{10}^{n+3}} +
\cdots +
\frac{d_{m}}{{10}^m}
\leq
\frac{1}{{10}^{n}}
\bigl(
1-{(\nicefrac{1}{10})}^{m-n}
\bigr)
<
\frac{1}{{10}^{n}} .
\end{equation*}
Dengan mengambil supremum terhadap $m$, kita peroleh
\begin{equation*}
x - D_n
\leq
\frac{1}{{10}^{n}} .
\end{equation*}

Selanjutnya kita beralih ke
butir kedua. Ambil sebarang $x \in (0,1]$.
Pertama, kita membuktikan eksistensinya.
Untuk memudahkan, tetapkan $D_0 \coloneqq 0$.
Maka,
$D_0 < x \leq D_0 + {10}^{-0}$.
Misalkan kita telah mendefinisikan digit-digit $d_1,d_2,\ldots,d_n$,
dan bahwa
$D_k < x \leq D_k + {10}^{-k}$, untuk $k=0,1,2,\ldots,n$. Kita perlu mendefinisikan $d_{n+1}$.

Berdasarkan
\hyperref[thm:arch:i]{sifat Archimedes} bilangan real,
terdapat bilangan bulat $j$ sedemikian sehingga
$x-D_n \leq j {10}^{-(n+1)}$. Ambil $j$ terkecil yang memenuhi sifat tersebut dan diperoleh
\begin{equation} \label{eq:theDnjineq}
(j-1){10}^{-(n+1)} < x-D_n \leq j {10}^{-(n+1)} .
\end{equation}
Tetapkan $d_{n+1} \coloneqq j-1$.
Karena $D_n < x$,
kita memiliki $d_{n+1} = j-1 \geq 0$. Di sisi lain, karena
$x-D_n \leq {10}^{-n}$, kita memperoleh bahwa
$j$ paling banyak 10, dan oleh karena itu $d_{n+1} \leq 9$.
Jadi, $d_{n+1}$ merupakan
digit desimal.
Karena $D_{n+1} = D_n + d_{n+1} {10}^{-(n+1)}$,
tambahkan $D_n$ pada ketaksamaan
\eqref{eq:theDnjineq} di atas:
\begin{multline*}
D_{n+1} = D_n + (j-1){10}^{-(n+1)} < x
\leq
D_n + j {10}^{-(n+1)} \\
=
D_n + (j-1) {10}^{-(n+1)} +
{10}^{-(n+1)} = D_{n+1} + {10}^{-(n+1)} .
\end{multline*}
Dengan demikian,
$D_{n+1} < x \leq D_{n+1} + {10}^{-(n+1)}$ berlaku.
Secara induktif, kita telah
mendefinisikan barisan digit tak hingga $0.d_1d_2d_3\ldots$.

Tinjau $D_{n} < x \leq D_{n} + {10}^{-n}$.
Karena $D_n < x$ untuk semua $n$, kita memiliki
$\sup \{ D_n : n \in \N \} \leq x$.
Ketaksamaan kedua untuk $D_n$ mengakibatkan
\begin{equation*}
x - \sup \{ D_m : m \in \N \}
\leq
x - D_n \leq 10^{-n} .
\end{equation*}
Karena ketaksamaan tersebut berlaku untuk semua $n$ dan
${10}^{-n}$ dapat dibuat sekecil apa pun
(lihat \exerciseref{exercise:bnlimit}),
kita memiliki $x \leq \sup \{ D_m : m \in \N \}$.
Oleh karena itu, $\sup \{ D_m : m \in \N \} = x$.

Yang tersisa untuk dibuktikan adalah keunikannya.
Misalkan $0.e_1e_2e_3\ldots$ merupakan representasi lain dari $x$.
Misalkan $E_n$ adalah hasil pemotongan hingga $n$ digit dari $0.e_1e_2e_3\ldots$, dan andaikan
$E_n < x \leq E_n + {10}^{-n}$ untuk semua $n \in \N$.
Misalkan untuk suatu $K \in \N$, $e_n = d_n$ untuk semua $n < K$, sehingga
$D_{K-1} = E_{K-1}$. Maka
\begin{equation*}
E_K = D_{K-1} + e_K{10}^{-K} < x \leq E_K + {10}^{-K} = D_{K-1} +
e_K{10}^{-K} + {10}^{-K} .
\end{equation*}
Dengan mengurangkan $D_{K-1}$ dan mengalikan dengan ${10}^{K}$, kita memperoleh
\begin{equation*}
e_K < (x - D_{K-1}){10}^K \leq e_K + 1 .
\end{equation*}
Dengan cara serupa,
\begin{equation*}
d_K < (x - D_{K-1}){10}^K \leq d_K + 1 .
\end{equation*}
Jadi, baik $e_K$ maupun $d_K$ merupakan bilangan bulat terbesar $j$
sedemikian sehingga $j < (x - D_{K-1}){10}^K$, dan oleh karena itu $e_K = d_K$.
Dengan kata lain, representasi tersebut unik.
\end{proof}

Representasi tersebut tidak unik jika kita tidak mensyaratkan
$D_n < x$ untuk semua $n$.
Sebagai contoh, untuk
bilangan $\nicefrac{1}{2}$, metode dalam pembuktian menghasilkan representasi
\begin{equation*}
0.49999\ldots .
\end{equation*}
Namun, $\nicefrac{1}{2}$ juga memiliki representasi $0.50000\ldots$.

Satu-satunya bilangan yang memiliki representasi tidak unik
adalah bilangan yang berakhir dengan barisan tak hingga digit $0$
atau barisan tak hingga digit $9$, karena satu-satunya representasi yang memenuhi
$D_n = x$ adalah representasi dengan semua digit setelah digit ke-$n$ bernilai nol.
Dalam kasus ini,
terdapat tepat dua representasi dari $x$ (lihat latihan).

Mari kita berikan pembuktian lain bahwa himpunan bilangan real tak terhitung dengan
menggunakan representasi desimal.
Ini adalah pembuktian kedua Cantor, yang mungkin lebih dikenal.
Pembuktian ini mungkin tampak lebih singkat, tetapi itu karena kita telah menyelesaikan
bagian sulitnya di atas dan kini hanya tersisa sebuah trik cerdik untuk membuktikan bahwa $\R$
tak terhitung.  Trik ini disebut
\emph{\myindex{diagonalisasi Cantor}}\index{diagonalisasi} dan
juga digunakan dalam pembuktian-pembuktian lain.

\begin{thm}[Cantor]
Himpunan $(0,1]$ tak terhitung.
\end{thm}

\begin{proof}
Misalkan $X \coloneqq \{ x_1,x_2,x_3,\ldots \}$ adalah sembarang himpunan bagian terhitung dari bilangan real dalam $(0,1]$.
Kita akan mengonstruksi suatu bilangan real yang tidak berada di $X$.  Misalkan
\begin{equation*}
x_n = 0.d_1^nd_2^nd_3^n\ldots
\end{equation*}
merupakan representasi unik dari proposisi tersebut; artinya, $d_j^n$ adalah
digit ke-$j$ dari bilangan ke-$n$.  Misalkan
\begin{equation*}
e_n \coloneqq
\begin{cases}
1 & \text{jika } d_n^n \neq 1, \\
2 & \text{jika } d_n^n = 1.
\end{cases}
\end{equation*}
Misalkan $E_n$ adalah hasil pemotongan hingga $n$ digit dari $y = 0.e_1e_2e_3\ldots$.  Karena
semua digitnya bukan nol, $E_n < E_{n+1} \leq y$.  Oleh karena itu,
\begin{equation*}
E_n < y \leq E_n + {10}^{-n}
\end{equation*}
untuk semua $n$, dan representasi tersebut adalah representasi unik untuk $y$ dari
proposisi tersebut.  Untuk setiap $n$, digit ke-$n$
dari $y$ berbeda dengan digit ke-$n$ dari $x_n$, sehingga $y \neq x_n$.
Jadi $y \notin X$, dan karena $X$ merupakan sembarang himpunan bagian terhitung,
$(0,1]$ harus tak terhitung.  Lihat \figureref{cantorexamplefig} untuk
contohnya.
\begin{myfigureht}
\myincludepdft{cantorexamplefig}{%
Lima bilangan bertanda x subskrip 1 hingga x subskrip 5 diperlihatkan tersusun sejajar, satu di atas
yang lain, dalam notasi desimal.
Lima digit desimal pertama dari setiap bilangan diperlihatkan.
Secara khusus, x subskrip 1 adalah 0 koma 1 3 2 1 0 dan elipsis yang menyiratkan
bahwa digit-digitnya berlanjut.
Digit pertama, yaitu 1, ditandai dengan sebuah kotak.
x subskrip 2 adalah 0 koma 7 9 4 1 3 dan elipsis.
Digit kedua, yaitu 9, ditandai dengan sebuah kotak.
x subskrip 3 adalah 0 koma 3 0 1 3 4 dan elipsis.
Digit ketiga, yaitu 1, ditandai dengan sebuah kotak.
x subskrip 4 adalah 0 koma 8 9 2 5 6 dan elipsis.
Digit keempat, yaitu 5, ditandai dengan sebuah kotak.
x subskrip 5 adalah 0 koma 1 6 0 2 4 dan elipsis.
Digit kelima, yaitu 4, ditandai dengan sebuah kotak.
Elipsis selanjutnya menunjukkan bahwa barisan bilangan terus berlanjut.
Di sebelah kanan terdapat teks: Bilangan yang tidak ada dalam daftar:
y sama dengan 0 koma 2 1 2 1 1 elipsis.}
\caption{Contoh diagonalisasi Cantor, dengan digit diagonal $d_n^n$
ditandai.\label{cantorexamplefig}}
\end{myfigureht}
\end{proof}

Dengan digit desimal, kita juga dapat menemukan banyak bilangan irasional.
Proposisi berikut berlaku untuk setiap
bilangan rasional, tetapi demi kesederhanaan kita hanya menyatakannya untuk $x \in (0,1]$.

\begin{prop} \label{prop:rationaldecimal}
Jika $x \in (0,1]$ adalah bilangan rasional dan $x = 0.d_1d_2d_3\ldots$,
maka digit desimalnya akhirnya mulai berulang.  Artinya, terdapat
bilangan bulat positif $N$ dan $P$, sedemikian sehingga untuk semua $n \geq N$, $d_n = d_{n+P}$.
\end{prop}

\begin{proof}
Misalkan $x = \nicefrac{p}{q}$ untuk bilangan bulat positif $p$ dan $q$.
Jika $x$ memiliki representasi desimal yang tidak unik, kedua representasinya
pada akhirnya memiliki digit berulang,
baik semuanya $0$ maupun semuanya $9$;
lihat \exerciseref{exercise:nonuniquedecimals}.
Oleh karena itu, misalkan $x$ memiliki representasi unik.
Kita juga misalkan bahwa $x \neq 1$ sehingga $p < q$.

Untuk menghitung digit pertama, ambil $10 p$ dan bagi dengan
$q$.  Misalkan $d_1$ adalah hasil baginya.
Sisa $r_1$ adalah suatu bilangan bulat
antara $0$ dan $q-1$.  Artinya, $d_1$ adalah bilangan bulat terbesar
sedemikian sehingga $d_1 q \leq 10p$, lalu $r_1 = 10p - d_1q$.
Karena $p < q$, kita memperoleh $d_1 < 10$, sehingga $d_1$ adalah sebuah digit.
Selain itu,
\begin{equation*}
\frac{d_1}{10} \leq \frac{p}{q} =
\frac{d_1}{10} + \frac{r_1}{10q} \leq \frac{d_1}{10} +
\frac{1}{10} .
\end{equation*}
Ketaksamaan pertama bersifat ketat karena $x$ memiliki
representasi unik.  Jadi $d_1$ memang merupakan digit pertama.
Bagian yang tersisa adalah $\frac{r_1}{10q}$.
Untuk menghitung $d_2$, bagi $10 r_1$ dengan $q$ sehingga diperoleh $d_2$
dan $r_2$, dengan $r_2 = 10 r_1 - d_2 q$.
Nyatakan $r_1$ lalu substitusikan ke dalam
$\frac{p}{q} = \frac{d_1}{10} + \frac{r_1}{10 q}$
sehingga diperoleh
$\frac{p}{q} = \frac{d_1}{10} + \frac{d_2}{{10}^2} +
\frac{r_2}{{10}^2 q} = D_2 +
\frac{r_2}{{10}^2 q}$.
Karena $r_2 < q$, sekali lagi kita memperoleh $d_2 < 10$.
Untuk menghitung $d_3$, bagi $10 r_2$ dengan $q$
sehingga diperoleh $d_3$ dan $r_3$.  Demikian seterusnya.
Setelah menghitung $n-1$ digit, kita memperoleh
$\frac{p}{q} = D_{n-1} + \frac{r_{n-1}}{10^{n-1} q}$.
Untuk mendapatkan digit ke-$n$,
bagi $10 r_{n-1}$ dengan $q$
untuk memperoleh hasil bagi $d_n$, sisa $r_n$, dan ketaksamaan
\begin{equation*}
\frac{d_n}{10} \leq \frac{r_{n-1}}{q} =
\frac{d_n}{10} + \frac{r_n}{10q} \leq \frac{d_n}{10} +
\frac{1}{10} .
\end{equation*}
Dengan membaginya dengan $10^{n-1}$ dan menambahkan $D_{n-1}$, kita memperoleh
\begin{equation*}
D_n \leq D_{n-1} + \frac{r_{n-1}}{10^{n-1} q} = \frac{p}{q} \leq D_n +
\frac{1}{10^n} .
\end{equation*}
Berdasarkan keunikan, kita memang memperoleh
digit ke-$n$, yaitu $d_n$, dari konstruksi ini.

Digit baru hanya bergantung pada sisa dari langkah
sebelumnya.
Terdapat paling banyak $q$ kemungkinan sisa.
Oleh karena itu, proses tersebut harus mulai berulang setelah paling banyak $q$ langkah,
sehingga $P$ paling banyak $q$.
\end{proof}

Kebalikan dari proposisi tersebut juga benar dan diserahkan sebagai latihan.

\begin{example}
Bilangan
\begin{equation*}
x = 0.101001000100001000001\ldots
\end{equation*}
adalah bilangan irasional.
Secara khusus, digit-digitnya adalah satu,
kemudian nol, kemudian satu, lalu dua nol, lalu satu,
kemudian tiga nol,
dan seterusnya.
Keirasionalan $x$ diperoleh dari
proposisi tersebut; digit-digitnya tidak pernah mulai berulang.
Untuk setiap $P$,
jika kita melangkah cukup jauh, kita menemukan digit 1 yang diikuti sekurang-kurangnya $P$ nol,
dan terdapat tak hingga banyaknya digit 1 semacam itu.
Artinya, terdapat tak hingga banyaknya $n$ (jadi $n$ dapat sebesar apa pun)
sedemikian sehingga $d_n=1$ dan $d_{n+P}=0$, dengan kata lain,
$d_n \neq d_{n+P}$.
\end{example}

\subsection{Latihan}

\begin{exercise}[Mudah]
Apa representasi desimal dari $1$ yang dijamin oleh
\propref{prop:decimalprop}?  Pastikan Anda menunjukkan bahwa representasi tersebut memang memenuhi
syaratnya.
\end{exercise}

\begin{exercise}
Buktikan kebalikan dari \propref{prop:rationaldecimal}, yaitu,
jika digit-digit dalam representasi desimal $x$ akhirnya berulang, maka
$x$ harus rasional.
\end{exercise}

\begin{exercise} \label{exercise:nonuniquedecimals}
Tunjukkan bahwa bilangan real $x \in (0,1)$ yang memiliki representasi desimal tidak unik
tepat merupakan bilangan rasional yang dapat dituliskan
sebagai $\frac{m}{10^n}$ untuk bilangan bulat $m$ dan $n$ tertentu.  Dalam hal ini, tunjukkan bahwa
terdapat tepat dua representasi dari $x$.
\end{exercise}

\begin{exercise} \label{exercise:decimalpropbaseb}
Misalkan $b \geq 2$ merupakan bilangan bulat.  Definisikan representasi bilangan real dalam
$[0,1]$ dengan basis $b$, bukan basis 10, dan buktikan
\propref{prop:decimalprop} untuk basis $b$.
\end{exercise}

\begin{exercise}
Dengan menggunakan latihan sebelumnya untuk $b=2$ (biner),
tunjukkan bahwa kardinalitas $\R$ sama dengan kardinalitas $\sP(\N)$,
sehingga diperoleh satu lagi bukti (meskipun terkait) bahwa $\R$ tak terhitung.
Petunjuk: Konstruksikan dua injeksi, satu dari $[0,1]$ ke $\sP(\N)$
dan satu dari $\sP(\N)$ ke $[0,1]$.  Petunjuk 2: Diberikan suatu
himpunan $A \subset \N$, pemetaan yang menetapkan digit biner ke-$n$ dari $x$ sama dengan 1 jika
$n\in A$ hampir berhasil, tetapi tidak injektif.  Modifikasilah pemetaan itu.
\end{exercise}

\begin{exercise}[Menantang]
Konstruksikan secara eksplisit suatu injeksi dari $[0,1] \times [0,1]$ ke $[0,1]$ (pikirkan mengapa
hal ini begitu mengejutkan\footnote{%
Dengan usaha yang jauh lebih besar (atau dengan menerapkan teorema
Cantor--Bernstein--Schr\"oder),
dapat dibuktikan bahwa terdapat suatu bijeksi.
Ketika membuktikan hasil ini, Cantor konon menulis,
\myquote{Saya melihatnya, tetapi saya tidak memercayainya.}}).
Kemudian deskripsikan himpunan bilangan dalam $[0,1]$ yang tidak berada dalam citra
injeksi Anda (kecuali, tentu saja, Anda berhasil mengonstruksi suatu bijeksi).
Petunjuk: Tinjau digit pada posisi genap dan ganjil dalam ekspansi desimal.
\end{exercise}

\begin{exercise}
Buktikan bahwa jika $x = \nicefrac{p}{q} \in (0,1]$ merupakan bilangan rasional, $q > 1$,
maka periode $P$ dari digit berulang dalam representasi desimal
$x$ sebenarnya kurang dari atau sama dengan $q-1$.
\end{exercise}

\begin{exercise} \label{exercise:bnlimit}
Buktikan bahwa jika $b \in \N$ dan $b \geq 2$, maka untuk setiap $\epsilon > 0$,
terdapat $n \in \N$ sedemikian sehingga
$b^{-n} < \epsilon$.  Petunjuk:
Salah satu caranya adalah terlebih dahulu membuktikan bahwa $b^n > n$ untuk setiap $n \in \N$ dengan induksi.
\end{exercise}

\begin{exercise}
Konstruksikan secara eksplisit suatu injeksi $f \colon \R \to \R \setminus \Q$ dengan menggunakan
\propref{prop:rationaldecimal}.
\end{exercise}
